\chapter{Chapter04 — Hastad Switching}
%%% generated by the blueprint pipeline from FABL.Chapter04.HastadSwitching
\section{Overview}

Declarations of the Lean module \texttt{FABL.Chapter04.HastadSwitching}.

\section{Declarations}

\begin{definition}[Free literal positions]
\lean{FABL.DNFTerm.freeLiteralPositions}
\statementsource{arora-barak-ch13-circuit-lowerbounds}{
  pdf-52929dbaa962-p002-b002,
  pdf-52929dbaa962-p004-b001
}
\statementsource{Analysis-of-Boolean-Functions-by-Ryan-ODonnell}{
  pdf-ffbe9c7b3306-p093-b004,
  pdf-ffbe9c7b3306-p098-b004
}
\label{FABL.DNFTerm.freeLiteralPositions}
\leanok
\uses{FABL.CoordRestriction, FABL.DNFTerm}
Let $T$ be a DNF term on $n$ coordinates, and let $\rho$ assign to each coordinate one
of the statuses free, fixed to $1$, or fixed to $-1$. Define the free literal positions
to be the positions $j$ in the list of literals of $T$ for which the coordinate of the
$j$th literal is assigned status free by $\rho$.
\end{definition}

\begin{theorem}[Free literal positions]
\lean{FABL.DNFTerm.mem_freeLiteralPositions_iff}
\label{FABL.DNFTerm.mem_freeLiteralPositions_iff}
\leanok
\uses{FABL.CoordRestriction, FABL.DNFTerm, FABL.DNFTerm.freeLiteralPositions}
\difficulty{1}
Let $n$ be a natural number, let $T$ be a DNF term on $n$ coordinates, and let each
coordinate be assigned one of the restrictions free, fixed to $1$, or fixed to $-1$. A
valid position $j$ in the literal list belongs to the set of free literal positions
exactly when the coordinate of the literal at $j$ is assigned the free restriction.
\end{theorem}

\begin{theorem}[Free literal coordinates lie in the support]
\lean{FABL.DNFTerm.freeLiteralPosition_coordinate_mem_support}
\label{FABL.DNFTerm.freeLiteralPosition_coordinate_mem_support}
\leanok
\uses{FABL.CoordRestriction, FABL.DNFTerm, FABL.DNFTerm.freeLiteralPositions, FABL.DNFTerm.mem_support_of_mem_literals, FABL.DNFTerm.support}
\difficulty{1}
Let $T$ be a conjunction of literals on $n$ coordinates and let $\rho$ assign each
coordinate a restriction state. If position $j$ in the list of literals of $T$ is free
under $\rho$, then the coordinate of the literal at position $j$ belongs to the support
of $T$.
\end{theorem}

\begin{theorem}[Evaluation under a compatible coordinate restriction]
\lean{FABL.DNFTerm.eval_complete_eq_neg_one_of_compatible_of_no_free}
\label{FABL.DNFTerm.eval_complete_eq_neg_one_of_compatible_of_no_free}
\leanok
\uses{FABL.CoordRestriction, FABL.CoordRestriction.complete, FABL.DNFTerm, FABL.DNFTerm.Compatible, FABL.DNFTerm.eval, FABL.DNFTerm.eval_eq_neg_one_iff, FABL.DNFTerm.freeLiteralPositions, FABL.signEncode}
\difficulty{5}
Let $n$ be a nonnegative integer, and let $T$ be a conjunction of sign literals on the
$n$ coordinates, with no coordinate appearing more than once. Let $\rho$ assign each
coordinate either the status free, fixed to $1$, or fixed to $-1$, and let a completion
of $\rho$ be any sign vector that agrees with its fixed coordinates and assigns
arbitrary signs to its free coordinates. If some completion makes every literal of $T$
match its required sign, and no literal of $T$ lies at a free coordinate, then every
completion of $\rho$ makes every literal of $T$ match its required sign.
\end{theorem}

\begin{definition}[Free coordinates lie within the available set]
\lean{FABL.CoordRestriction.FreeWithin}
\statementsource{lmn-constant-depth-fourier}{pdf-63759bd909d5-p005-b001}
\statementsource{arora-barak-ch13-circuit-lowerbounds}{
  pdf-52929dbaa962-p002-b001,
  pdf-52929dbaa962-p002-b003
}
\statementsource{Analysis-of-Boolean-Functions-by-Ryan-ODonnell}{
  pdf-ffbe9c7b3306-p098-b003,
  pdf-ffbe9c7b3306-p098-b004
}
\label{FABL.CoordRestriction.FreeWithin}
\leanok
\uses{FABL.CoordRestriction}
Let $n$ be a natural number, let $\rho$ assign to each coordinate in $\mathrm{Fin}(n)$
one of the statuses free, fixed to $1$, or fixed to $-1$, and let
$A\subseteq\mathrm{Fin}(n)$. Say that $\rho$ is free within $A$ if every coordinate
assigned status free belongs to $A$.
\end{definition}

\begin{theorem}[Free coordinates lie in the full coordinate set]
\lean{FABL.CoordRestriction.freeWithin_univ}
\label{FABL.CoordRestriction.freeWithin_univ}
\leanok
\uses{FABL.CoordRestriction, FABL.CoordRestriction.FreeWithin}
\difficulty{3}
Let $n$ be a natural number and let $\rho$ assign a restriction state to each of the $n$
coordinates. Every coordinate that is free under $\rho$ belongs to the full set of
coordinates.
\end{theorem}

\begin{theorem}[A fixed sign gives a nonfree state]
\lean{FABL.CoordRestriction.fixedState_ne_free}
\label{FABL.CoordRestriction.fixedState_ne_free}
\leanok
\uses{FABL.CoordRestriction.fixedState, FABL.Sign}
\difficulty{4}
For every sign $s\in\{\pm1\}$, the restriction state determined by $s$ is not free.
\end{theorem}

\begin{theorem}[Free coordinates after fixing one coordinate]
\lean{FABL.CoordRestriction.freeWithin_fixCoordinate}
\label{FABL.CoordRestriction.freeWithin_fixCoordinate}
\leanok
\uses{FABL.CoordRestriction, FABL.CoordRestriction.FreeWithin, FABL.CoordRestriction.fixCoordinate, FABL.CoordRestriction.fixCoordinate_apply_of_ne, FABL.CoordRestriction.fixedState_ne_free, FABL.Sign}
\difficulty{4}
Let $n$ be a natural number, let $\rho$ assign a restriction state to each of the $n$
coordinates, and let $A$ be a set of coordinates such that every coordinate free under
$\rho$ belongs to $A$. For any coordinate $i$ and any sign $s\in\{1,-1\}$, after fixing
coordinate $i$ according to $s$, every coordinate that remains free belongs to
$A\setminus\{i\}$.
\end{theorem}

\begin{definition}[Fixing a List of Coordinates]
\lean{FABL.CoordRestriction.fixCoordinateList}
\statementsource{arora-barak-ch13-circuit-lowerbounds}{
  pdf-52929dbaa962-p002-b001,
  pdf-52929dbaa962-p002-b003
}
\statementsource{Analysis-of-Boolean-Functions-by-Ryan-ODonnell}{
  pdf-ffbe9c7b3306-p098-b003,
  pdf-ffbe9c7b3306-p098-b004
}
\label{FABL.CoordRestriction.fixCoordinateList}
\leanok
\uses{FABL.CoordRestriction, FABL.CoordRestriction.fixCoordinate, FABL.signEncode}
Let $\rho$ assign each of the $n$ coordinates a state that is either free, fixed at $1$,
or fixed at $-1$, and let $x\in\mathbb{F}_2^n$. Starting from $\rho$, process the
coordinates in the given list order, fixing each listed coordinate $i$ to the sign
encoded by $x_i$; the resulting assignment of states is the fixed-coordinate list.
\end{definition}

\begin{theorem}[Empty list of coordinate updates]
\lean{FABL.CoordRestriction.fixCoordinateList_nil}
\label{FABL.CoordRestriction.fixCoordinateList_nil}
\leanok
\uses{FABL.CoordRestriction, FABL.CoordRestriction.fixCoordinateList}
\difficulty{1}
Let $n$ be a natural number, let $\rho$ assign one of the three restriction states to
each of the $n$ coordinates, and let $x\in(\mathbb{Z}/2\mathbb{Z})^n$. Applying the
coordinate-fixing procedure to the empty list leaves $\rho$ unchanged.
\end{theorem}

\begin{theorem}[Fixing a list by its first coordinate]
\lean{FABL.CoordRestriction.fixCoordinateList_cons}
\label{FABL.CoordRestriction.fixCoordinateList_cons}
\leanok
\uses{FABL.CoordRestriction, FABL.CoordRestriction.fixCoordinate, FABL.CoordRestriction.fixCoordinateList, FABL.signEncode}
\difficulty{1}
Let $n$ be a natural number, let $\rho$ assign a restriction state to each of the $n$
coordinates, let $i$ be a coordinate, let $L$ be a list of coordinates, and let
$x\in\mathbb{F}_2^n$. Fixing the coordinates in the list consisting of $i$ followed by
$L$ according to $x$ gives the same assignment as first fixing $i$ according to the sign
encoded by $x_i$ and then fixing the coordinates in $L$ according to $x$.
\end{theorem}

\begin{theorem}[Completion after fixing a coordinate]
\lean{FABL.CoordRestriction.complete_fixCoordinate_signEncode}
\label{FABL.CoordRestriction.complete_fixCoordinate_signEncode}
\leanok
\uses{FABL.CoordRestriction, FABL.CoordRestriction.complete, FABL.CoordRestriction.complete_fixCoordinate_eq_of_value, FABL.CoordRestriction.fixCoordinate, FABL.signEncode}
\difficulty{3}
Let $n$ be a natural number, let $\rho$ assign a restriction state to each of the $n$
coordinates, let $i$ be one of those coordinates, and let $x\in\mathbb{F}_2^n$. If
coordinate $i$ is free in $\rho$, then completing $x$ after fixing coordinate $i$ to the
sign encoded by $x_i$ gives the same sign vector as completing $x$ using $\rho$.
\end{theorem}

\begin{theorem}[Completion after fixing a coordinate list]
\lean{FABL.CoordRestriction.complete_fixCoordinateList}
\label{FABL.CoordRestriction.complete_fixCoordinateList}
\leanok
\uses{FABL.CoordRestriction, FABL.CoordRestriction.complete, FABL.CoordRestriction.complete_fixCoordinate_signEncode, FABL.CoordRestriction.fixCoordinateList, FABL.CoordRestriction.fixCoordinateList_cons, FABL.CoordRestriction.fixCoordinate_apply_of_ne, FABL.signEncode}
\difficulty{4}
Let $n$ be a natural number, let $\rho$ be a coordinate restriction on $n$ coordinates,
let $L$ be a list of distinct coordinates, and let $x$ be a binary vector of length $n$.
Suppose every coordinate in $L$ is free in $\rho$. Completing the restriction obtained
by fixing the coordinates in $L$ according to the signs encoded by $x$, using $x$, gives
the same sign vector as completing $\rho$ using $x$.
\end{theorem}

\begin{theorem}[Value outside a listed coordinate set]
\lean{FABL.CoordRestriction.fixCoordinateList_apply_of_not_mem}
\label{FABL.CoordRestriction.fixCoordinateList_apply_of_not_mem}
\leanok
\uses{FABL.CoordRestriction, FABL.CoordRestriction.fixCoordinateList, FABL.CoordRestriction.fixCoordinateList_cons, FABL.CoordRestriction.fixCoordinate_apply_of_ne, FABL.signEncode}
\difficulty{4}
Let $n$ be a natural number, let $\rho$ assign a restriction state to each of the $n$
coordinates, let $L$ be a list of coordinates, and let $x\in\mathbb{F}_2^n$. At any
coordinate $j$ not in $L$, fixing the coordinates in $L$ according to the signs encoded
by $x$ leaves the state equal to $\rho_j$.
\end{theorem}

\begin{theorem}[Listed coordinates are fixed]
\lean{FABL.CoordRestriction.fixCoordinateList_ne_free_of_mem}
\label{FABL.CoordRestriction.fixCoordinateList_ne_free_of_mem}
\leanok
\uses{FABL.CoordRestriction, FABL.CoordRestriction.fixCoordinateList, FABL.CoordRestriction.fixCoordinateList_apply_of_not_mem, FABL.CoordRestriction.fixCoordinateList_cons, FABL.CoordRestriction.fixCoordinate_apply_self, FABL.CoordRestriction.fixedState_ne_free}
\difficulty{5}
Let $n$ be a natural number, let $\rho$ assign one of the three restriction states to
each of the $n$ coordinates, let $L$ be a list of coordinates, and let
$x\in(\mathbb{Z}/2\mathbb{Z})^n$. Every coordinate belonging to $L$ is nonfree after
successively fixing the listed coordinates according to the signs encoded by $x$.
\end{theorem}

\begin{theorem}[Free coordinates remain free]
\lean{FABL.CoordRestriction.fixCoordinateList_free_imp_original_free}
\label{FABL.CoordRestriction.fixCoordinateList_free_imp_original_free}
\leanok
\uses{FABL.CoordRestriction, FABL.CoordRestriction.fixCoordinateList, FABL.CoordRestriction.fixCoordinateList_apply_of_not_mem, FABL.CoordRestriction.fixCoordinateList_ne_free_of_mem}
\difficulty{4}
Let $n$ be a natural number, let $\rho$ assign one of the three restriction states to
each of the $n$ coordinates, let $L$ be a list of coordinates, and let
$x\in(\mathbb{Z}/2\mathbb{Z})^n$. If coordinate $j$ is free after successively fixing
the coordinates in $L$ according to the signs encoded by $x$, then $j$ was free in
$\rho$.
\end{theorem}

\begin{theorem}[Free coordinates after fixing a list]
\lean{FABL.CoordRestriction.freeWithin_fixCoordinateList}
\label{FABL.CoordRestriction.freeWithin_fixCoordinateList}
\leanok
\uses{FABL.CoordRestriction, FABL.CoordRestriction.FreeWithin, FABL.CoordRestriction.fixCoordinateList, FABL.CoordRestriction.fixCoordinateList_apply_of_not_mem, FABL.CoordRestriction.fixCoordinateList_ne_free_of_mem}
\difficulty{4}
Let $n$ be a natural number, let $\rho$ assign a restriction state to each of the $n$
coordinates, and let $A$ be a finite set of coordinates such that every coordinate free
under $\rho$ belongs to $A$. Let $L$ be a list of coordinates and let $x$ be a bit
assignment. After fixing the coordinates in $L$ according to $x$, every coordinate that
remains free belongs to $A\setminus L$, where $L$ is viewed as the set of its entries.
\end{theorem}

\begin{theorem}[Free coordinates after fixing a list]
\lean{FABL.CoordRestriction.card_freeCoordinates_fixCoordinateList_lt}
\label{FABL.CoordRestriction.card_freeCoordinates_fixCoordinateList_lt}
\leanok
\uses{FABL.CoordRestriction, FABL.CoordRestriction.fixCoordinateList, FABL.CoordRestriction.fixCoordinateList_apply_of_not_mem, FABL.CoordRestriction.fixCoordinateList_ne_free_of_mem, FABL.CoordRestriction.freeCoordinates, FABL.CoordRestriction.mem_freeCoordinates_iff}
\difficulty{4}
Let $n$ be a natural number, let $\rho$ be a coordinate restriction on $n$ coordinates,
let $L$ be a nonempty list of coordinates, and let $x$ be a binary vector of length $n$.
Suppose every coordinate in $L$ is free in $\rho$. After fixing the coordinates in $L$
according to the signs encoded by $x$, the number of free coordinates is strictly less
than it was in $\rho$.
\end{theorem}

\begin{definition}[Least compatible DNF term index]
\lean{FABL.DNFFormula.canonicalTermIndex}
\label{FABL.DNFFormula.canonicalTermIndex}
\leanok
\uses{FABL.CoordRestriction, FABL.DNFFormula, FABL.DNFFormula.compatibleTermIndices}
Let $\varphi$ be a DNF formula on $n$ sign coordinates, and let $\rho$ assign to each
coordinate either no restriction, the value $1$, or the value $-1$. If at least one term
of $\varphi$ is compatible with $\rho$, define the canonical index to be the least index
in the formula's term list of a compatible term.
\end{definition}

\begin{theorem}[Compatibility of the canonical term]
\lean{FABL.DNFFormula.canonicalTerm_compatible}
\label{FABL.DNFFormula.canonicalTerm_compatible}
\leanok
\uses{FABL.CoordRestriction, FABL.DNFFormula, FABL.DNFFormula.canonicalTermIndex, FABL.DNFFormula.compatibleTermIndices, FABL.DNFTerm.Compatible}
\difficulty{1}
Let $\phi$ be a DNF formula on $n$ coordinates and let $\rho$ be a coordinate
restriction. If at least one term of $\phi$ is compatible with $\rho$, then the term of
least index compatible with $\rho$ is compatible with $\rho$.
\end{theorem}

\begin{definition}[Coordinate decision tree]
\lean{FABL.CoordinateDecisionTree}
\label{FABL.CoordinateDecisionTree}
\leanok
A coordinate decision tree with coordinates in $\{1,\ldots,n\}$ and values in a type
$\alpha$ is either a leaf labeled by an element of $\alpha$, or a query at one
coordinate with a zero-child and a one-child, each of which is a coordinate decision
tree of the same kind.
\end{definition}

\begin{definition}[Evaluation of a coordinate decision tree]
\lean{FABL.CoordinateDecisionTree.eval}
\label{FABL.CoordinateDecisionTree.eval}
\leanok
\uses{FABL.CoordinateDecisionTree}
Let $T$ be a coordinate decision tree with values in $\alpha$, and let
$x\in\mathbb{F}_2^n$. Its evaluation at $x$ is the leaf label if $T$ is a leaf; at a
query of coordinate $i$, it is the evaluation of the zero-child when $x_i=0$, and
otherwise the evaluation of the one-child.
\end{definition}

\begin{definition}[Depth of a coordinate decision tree]
\lean{FABL.CoordinateDecisionTree.depth}
\label{FABL.CoordinateDecisionTree.depth}
\leanok
\uses{FABL.CoordinateDecisionTree}
The depth of a coordinate decision tree is zero for a leaf and, for a query node, one
plus the maximum of the depths of its two child trees.
\end{definition}

\begin{definition}[Well-formed coordinate decision tree]
\lean{FABL.CoordinateDecisionTree.WellFormed}
\label{FABL.CoordinateDecisionTree.WellFormed}
\leanok
\uses{FABL.CoordinateDecisionTree}
Let $T$ be a coordinate decision tree and let $A$ be a finite set of available
coordinates. The tree is well formed relative to $A$ if it is a leaf, or if its query
coordinate belongs to $A$ and both child trees are well formed relative to the set
obtained by removing that coordinate from $A$.
\end{definition}

\begin{definition}[Conversion to a restricted decision tree]
\lean{FABL.CoordinateDecisionTree.toF₂DecisionTree}
\leanok
\uses{FABL.CoordinateDecisionTree, FABL.CoordinateDecisionTree.WellFormed, FABL.F₂DecisionTree}
Let $T$ be a coordinate decision tree with values in a type $\alpha$, and let $A$ be a
set of available coordinates among the $n$ coordinates. If every query in $T$ uses a
coordinate in the currently available set, and both children of each query satisfy the
same condition after that coordinate is removed, then $T$ determines a decision tree
indexed by $A$: leaves retain their values, and each query has the same two children,
indexed by the remaining coordinates.
\end{definition}

\begin{theorem}[Evaluation of a tree conversion]
\lean{FABL.CoordinateDecisionTree.eval_toF₂DecisionTree}
\leanok
\uses{FABL.CoordinateDecisionTree, FABL.CoordinateDecisionTree.WellFormed, FABL.CoordinateDecisionTree.eval, FABL.CoordinateDecisionTree.toF₂DecisionTree, FABL.F₂DecisionTree.eval}
\difficulty{4}
Let $T$ be a coordinate decision tree on the Boolean cube $\mathbb{F}_2^n$, with values
in a type $\alpha$, and let $A$ be a set of available coordinates. If every query made
by $T$ uses a coordinate in $A$ and each child tree satisfies the same condition after
that coordinate is removed, then converting $T$ to a decision tree whose queries
explicitly track available coordinates preserves its evaluation at every input in
$\mathbb{F}_2^n$.
\end{theorem}

\begin{theorem}[Depth under tree conversion]
\lean{FABL.CoordinateDecisionTree.depth_toF₂DecisionTree}
\leanok
\uses{FABL.CoordinateDecisionTree, FABL.CoordinateDecisionTree.WellFormed, FABL.CoordinateDecisionTree.depth, FABL.CoordinateDecisionTree.toF₂DecisionTree, FABL.F₂DecisionTree.depth}
\difficulty{4}
Let $n$ be a nonnegative integer, let $\alpha$ be a type, and let $T$ be a coordinate
decision tree with values in $\alpha$. For any set of available coordinates for which
$T$ is well formed—that is, every queried coordinate is available and both child trees
are well formed after that coordinate is removed—the depth of the converted tree, which
only queries available coordinates, equals the depth of $T$.
\end{theorem}

\begin{definition}[Assigning coordinates in a Boolean vector]
\lean{FABL.CoordinateDecisionTree.assignCoordinates}
\label{FABL.CoordinateDecisionTree.assignCoordinates}
\leanok
Let $s,x\in\mathbb F_2^n$ and let $i_1,\ldots,i_k$ be a list of coordinates. Starting
from $s$, successively replace each listed coordinate $i_j$ by its value in $x$, in list
order, and leave all other coordinates unchanged. The resulting vector is the assigned
vector.
\end{definition}

\begin{definition}[A tree querying a coordinate list]
\lean{FABL.CoordinateDecisionTree.queryCoordinates}
\label{FABL.CoordinateDecisionTree.queryCoordinates}
\leanok
\uses{FABL.CoordinateDecisionTree}
Let $i_1,\ldots,i_k$ be a list of coordinates among $n$, let $s\in\mathbb F_2^n$, and
let $F$ assign a coordinate vector to a coordinate decision tree with values in
$\alpha$. The recursively constructed tree queries the listed coordinates in order; at
each query, its zero child continues from the current vector with that coordinate set to
$0$, and its one child continues with that coordinate set to $1$. After the list is
exhausted, the tree is $F$ evaluated at the resulting vector.
\end{definition}

\begin{theorem}[Evaluation of a coordinate query tree]
\lean{FABL.CoordinateDecisionTree.eval_queryCoordinates}
\label{FABL.CoordinateDecisionTree.eval_queryCoordinates}
\leanok
\uses{FABL.CoordinateDecisionTree, FABL.CoordinateDecisionTree.assignCoordinates, FABL.CoordinateDecisionTree.eval, FABL.CoordinateDecisionTree.queryCoordinates}
\difficulty{4}
Let $n$ be a nonnegative integer, let $\alpha$ be a type, let $L$ be a list of
coordinates in $[n]$, and let $s,x$ be binary assignments on $[n]$. For any function
assigning a coordinate decision tree to each binary assignment, the tree that queries
the coordinates in $L$ and branches on whether each queried bit is $0$ or $1$, then uses
that function at the resulting assignment, evaluates at $x$ to the same value as the
tree assigned to the seed updated with the values of $x$ at the coordinates in $L$,
evaluated at $x$.
\end{theorem}

\begin{theorem}[Unqueried coordinates remain unchanged]
\lean{FABL.CoordinateDecisionTree.assignCoordinates_apply_of_not_mem}
\label{FABL.CoordinateDecisionTree.assignCoordinates_apply_of_not_mem}
\leanok
\uses{FABL.CoordinateDecisionTree.assignCoordinates}
\difficulty{4}
Let $n$ be a nonnegative integer, let $L$ be a list of coordinates in $	ext{Fin}(n)$,
and let $s,x$ be assignments of bits in $[2]$ to those coordinates. If $i$ is a
coordinate not in $L$, then after successively assigning to $s$ the values of $x$ at the
coordinates in $L$, the value at $i$ is still $s(i)$.
\end{theorem}

\begin{theorem}[Queried coordinates take input values]
\lean{FABL.CoordinateDecisionTree.assignCoordinates_eq_on_mem}
\label{FABL.CoordinateDecisionTree.assignCoordinates_eq_on_mem}
\leanok
\uses{FABL.CoordinateDecisionTree.assignCoordinates, FABL.CoordinateDecisionTree.assignCoordinates_apply_of_not_mem}
\difficulty{5}
Let $n$ be a nonnegative integer, let $L$ be a list of coordinates in $[n]$, and let
$s,x$ be binary assignments on $[n]$. After successively assigning from $x$ the values
at the coordinates in $L$ to the seed assignment $s$, every coordinate appearing in $L$
has its value from $x$.
\end{theorem}

\begin{theorem}[Restriction depends on listed coordinates]
\lean{FABL.CoordinateDecisionTree.fixCoordinateList_congr}
\label{FABL.CoordinateDecisionTree.fixCoordinateList_congr}
\leanok
\uses{FABL.CoordRestriction, FABL.CoordRestriction.fixCoordinateList, FABL.CoordRestriction.fixCoordinateList_cons}
\difficulty{4}
Let $\rho$ be a coordinate restriction, let $C$ be a list of coordinates, and let
$x,y\in\mathbb{F}_2^n$. If $x$ and $y$ agree at every coordinate appearing in $C$, then
fixing the coordinates in $C$ according to either vector produces the same restriction.
\end{theorem}

\begin{theorem}[Restriction after coordinate assignment]
\lean{FABL.CoordinateDecisionTree.fixCoordinateList_assignCoordinates}
\label{FABL.CoordinateDecisionTree.fixCoordinateList_assignCoordinates}
\leanok
\uses{FABL.CoordRestriction, FABL.CoordRestriction.fixCoordinateList, FABL.CoordinateDecisionTree.assignCoordinates, FABL.CoordinateDecisionTree.assignCoordinates_eq_on_mem, FABL.CoordinateDecisionTree.fixCoordinateList_congr}
\difficulty{3}
Let $\rho$ be a coordinate restriction, let $C$ be a list of coordinates, and let
$s,x\in\mathbb{F}_2^n$. Fixing the coordinates in $C$ according to the values of $x$
gives the same restriction whether those values are read directly from $x$ or from the
vector obtained by updating $s$ on $C$ with the corresponding values of $x$.
\end{theorem}

\begin{theorem}[Well-formedness of a queried-coordinate tree]
\lean{FABL.CoordinateDecisionTree.wellFormed_queryCoordinates}
\label{FABL.CoordinateDecisionTree.wellFormed_queryCoordinates}
\leanok
\uses{FABL.CoordinateDecisionTree, FABL.CoordinateDecisionTree.WellFormed, FABL.CoordinateDecisionTree.queryCoordinates}
\difficulty{5}
Let $C$ be a list of distinct coordinates, each belonging to a finite set $A$, let
$s\in\mathbb{F}_2^n$, and let $F$ assign a coordinate decision tree to each Boolean-cube
vector. If every tree $F(x)$ is well formed with respect to the coordinates in $A$
outside $C$, then the tree formed by querying the coordinates in $C$ from seed $s$, with
both bit branches at each query, is well formed with respect to $A$.
\end{theorem}

\begin{definition}[Coordinates at selected literal positions]
\lean{FABL.DNFTerm.coordinatesAtPositions}
\label{FABL.DNFTerm.coordinatesAtPositions}
\leanok
\uses{FABL.DNFTerm}
Let $T$ be a term on $n$ variables whose literals have pairwise distinct coordinates.
For any set $P$ of valid positions in its literal list, form the list of coordinates of
the literals at those positions, in the order induced by the set's increasing
enumeration.
\end{definition}

\begin{theorem}[Distinct coordinates at selected positions]
\lean{FABL.DNFTerm.coordinatesAtPositions_nodup}
\label{FABL.DNFTerm.coordinatesAtPositions_nodup}
\leanok
\uses{FABL.DNFTerm, FABL.DNFTerm.coordinatesAtPositions}
\difficulty{4}
Let $T$ be a conjunction of literals on $n$ coordinates whose literal indices are
pairwise distinct. For any set of positions in its list of literals, the coordinates at
those positions are pairwise distinct.
\end{theorem}

\begin{theorem}[Coordinates selected by positions]
\lean{FABL.DNFTerm.mem_coordinatesAtPositions_iff}
\label{FABL.DNFTerm.mem_coordinatesAtPositions_iff}
\leanok
\uses{FABL.DNFTerm, FABL.DNFTerm.coordinatesAtPositions}
\difficulty{1}
Let $n$ be a natural number, let $T$ be a DNF term on $n$ coordinates, and let $P$ be a
set of positions in its literal list. For any coordinate $i$, the coordinate $i$ occurs
among the literals at positions in $P$ if and only if some position in $P$ contains a
literal whose coordinate is $i$.
\end{theorem}

\begin{theorem}[Nonempty Selected Coordinates]
\lean{FABL.DNFTerm.coordinatesAtPositions_nonempty_iff}
\label{FABL.DNFTerm.coordinatesAtPositions_nonempty_iff}
\leanok
\uses{FABL.DNFTerm, FABL.DNFTerm.coordinatesAtPositions}
\difficulty{1}
Let $T$ be a conjunction of literals with pairwise distinct variable coordinates, and
let $P$ be a set of positions in its list of literals. The list of coordinates selected
at positions in $P$ is nonempty if and only if $P$ is nonempty.
\end{theorem}

\begin{definition}[Recursive canonical block tree]
\lean{FABL.DNFFormula.canonicalBlockTreeAux}
\label{FABL.DNFFormula.canonicalBlockTreeAux}
\leanok
\uses{FABL.CoordRestriction, FABL.CoordRestriction.complete, FABL.CoordRestriction.fixCoordinateList, FABL.CoordinateDecisionTree, FABL.CoordinateDecisionTree.queryCoordinates, FABL.DNFFormula, FABL.DNFFormula.canonicalTermIndex, FABL.DNFFormula.compatibleTermIndices, FABL.DNFFormula.eval, FABL.DNFTerm.Compatible, FABL.DNFTerm.coordinatesAtPositions, FABL.DNFTerm.freeLiteralPositions, FABL.Sign}
Let $\varphi$ be a disjunctive normal form on $n$ sign variables, and let $\rho$ assign
each coordinate one of three states: free, fixed to $1$, or fixed to $-1$. With fuel
$m$, define a coordinate decision tree recursively. At fuel zero, its leaf is the value
of $\varphi$ on the completion of $\rho$ by the zero vector. At positive fuel, if no
term is compatible with $\rho$, return the leaf $1$; otherwise select the compatible
term of least index. If it has no free literal positions, return the leaf $-1$. If it
has free literal positions, query their coordinates in order, beginning from the zero
vector; after those queries, fix the queried coordinates according to the resulting
vector, return the leaf $-1$ if the selected term is compatible with the updated
restriction, and otherwise continue recursively with one less unit of fuel and that
updated restriction.
\end{definition}

\begin{definition}[Raw canonical block tree]
\lean{FABL.DNFFormula.canonicalBlockTreeRaw}
\label{FABL.DNFFormula.canonicalBlockTreeRaw}
\leanok
\uses{FABL.CoordRestriction, FABL.CoordinateDecisionTree, FABL.DNFFormula, FABL.DNFFormula.canonicalBlockTreeAux, FABL.Sign}
Let $\varphi$ be a disjunctive normal form on $n$ sign variables and let $\rho$ assign
each coordinate a free or fixed sign state. The raw canonical block tree for $\varphi$
under $\rho$ is the recursive canonical block tree with fuel $n$.
\end{definition}

\begin{theorem}[Evaluation when no term is compatible]
\lean{FABL.DNFFormula.eval_eq_one_of_no_compatible}
\label{FABL.DNFFormula.eval_eq_one_of_no_compatible}
\leanok
\uses{FABL.CoordRestriction, FABL.CoordRestriction.complete, FABL.DNFFormula, FABL.DNFFormula.compatibleTermIndices, FABL.DNFFormula.eval, FABL.DNFFormula.eval_eq_neg_one_iff}
\difficulty{4}
Let $n$ be a nonnegative integer, let $\varphi$ be a DNF formula on $n$ sign
coordinates, and let $\rho$ be a coordinate restriction. If no term of $\varphi$ can
evaluate to $-1$ on any sign vector completing $\rho$, then for every $x\in\mathbb
F_2^n$, $\varphi$ evaluates to $1$ on the completion of $\rho$ specified by $x$.
\end{theorem}

\begin{theorem}[Completion with no free coordinates]
\lean{FABL.CoordRestriction.complete_eq_of_freeCoordinates_eq_empty}
\label{FABL.CoordRestriction.complete_eq_of_freeCoordinates_eq_empty}
\leanok
\uses{FABL.CoordRestriction, FABL.CoordRestriction.complete, FABL.CoordRestriction.freeCoordinates, FABL.CoordRestriction.mem_freeCoordinates_iff}
\difficulty{4}
Let $n$ be a natural number and let $\rho$ be a coordinate restriction on $n$
coordinates with no free coordinates. For any two binary vectors $x$ and $y$ of length
$n$, completing $\rho$ using $x$ gives the same sign vector as completing it using $y$.
\end{theorem}

\begin{theorem}[Coordinates at free literal positions]
\lean{FABL.DNFTerm.coordinatesAt_freeLiteralPositions_all_free}
\label{FABL.DNFTerm.coordinatesAt_freeLiteralPositions_all_free}
\leanok
\uses{FABL.CoordRestriction, FABL.DNFTerm, FABL.DNFTerm.coordinatesAtPositions, FABL.DNFTerm.freeLiteralPositions, FABL.DNFTerm.mem_coordinatesAtPositions_iff, FABL.DNFTerm.mem_freeLiteralPositions_iff}
\difficulty{4}
Let $n$ be a natural number, let $T$ be a conjunction of literals on $n$ coordinates
with no coordinate appearing more than once, and assign to each coordinate one of the
restrictions free, fixed at $+1$, or fixed at $-1$. Every coordinate belonging to a
literal whose position in $T$ has the free restriction also has the free restriction.
\end{theorem}

\begin{theorem}[Eliminating free literal positions]
\lean{FABL.DNFTerm.freeLiteralPositions_fixCoordinateList_eq_empty}
\label{FABL.DNFTerm.freeLiteralPositions_fixCoordinateList_eq_empty}
\leanok
\uses{FABL.CoordRestriction, FABL.CoordRestriction.fixCoordinateList, FABL.CoordRestriction.fixCoordinateList_free_imp_original_free, FABL.CoordRestriction.fixCoordinateList_ne_free_of_mem, FABL.DNFTerm, FABL.DNFTerm.coordinatesAtPositions, FABL.DNFTerm.freeLiteralPositions, FABL.DNFTerm.mem_coordinatesAtPositions_iff, FABL.DNFTerm.mem_freeLiteralPositions_iff}
\difficulty{4}
Let $T$ be a conjunction of literals on $n$ coordinates whose literal coordinates are
pairwise distinct, let $\rho$ assign each coordinate a free or fixed sign state, and let
$x\in\mathbb{F}_2^n$. Fix, according to the sign encoded by $x$, every coordinate that
is free and occurs in $T$. After these coordinates are fixed, no literal of $T$ remains
at a free coordinate.
\end{theorem}

\begin{theorem}[Evaluation of the canonical block tree]
\lean{FABL.DNFFormula.eval_canonicalBlockTreeAux}
\label{FABL.DNFFormula.eval_canonicalBlockTreeAux}
\leanok
\uses{FABL.CoordRestriction, FABL.CoordRestriction.card_freeCoordinates_fixCoordinateList_lt, FABL.CoordRestriction.complete, FABL.CoordRestriction.complete_eq_of_freeCoordinates_eq_empty, FABL.CoordRestriction.complete_fixCoordinateList, FABL.CoordRestriction.fixCoordinateList, FABL.CoordRestriction.freeCoordinates, FABL.CoordinateDecisionTree.assignCoordinates, FABL.CoordinateDecisionTree.eval, FABL.CoordinateDecisionTree.eval_queryCoordinates, FABL.CoordinateDecisionTree.fixCoordinateList_assignCoordinates, FABL.DNFFormula, FABL.DNFFormula.canonicalBlockTreeAux, FABL.DNFFormula.canonicalTermIndex, FABL.DNFFormula.canonicalTerm_compatible, FABL.DNFFormula.compatibleTermIndices, FABL.DNFFormula.eval, FABL.DNFFormula.eval_eq_neg_one_iff, FABL.DNFFormula.eval_eq_one_of_no_compatible, FABL.DNFTerm.Compatible, FABL.DNFTerm.coordinatesAtPositions, FABL.DNFTerm.coordinatesAtPositions_nodup, FABL.DNFTerm.coordinatesAtPositions_nonempty_iff, FABL.DNFTerm.coordinatesAt_freeLiteralPositions_all_free, FABL.DNFTerm.eval_complete_eq_neg_one_of_compatible_of_no_free, FABL.DNFTerm.freeLiteralPositions, FABL.DNFTerm.freeLiteralPositions_fixCoordinateList_eq_empty}
\difficulty{6}
Let $n$ be a nonnegative integer, let $\varphi$ be a DNF formula on $n$ sign
coordinates, let $\rho$ be a coordinate restriction, and let $k$ be a nonnegative
integer such that the number of free coordinates under $\rho$ is at most $k$. For every
$x\in\mathbb F_2^n$, the canonical block tree constructed with fuel $k$ evaluates at $x$
to the value of $\varphi$ on the sign vector obtained by completing $\rho$ using $x$.
\end{theorem}

\begin{theorem}[Well-formedness of the block tree construction]
\lean{FABL.DNFFormula.canonicalBlockTreeAux_wellFormed}
\label{FABL.DNFFormula.canonicalBlockTreeAux_wellFormed}
\leanok
\uses{FABL.CoordRestriction, FABL.CoordRestriction.FreeWithin, FABL.CoordRestriction.fixCoordinateList, FABL.CoordRestriction.freeWithin_fixCoordinateList, FABL.CoordinateDecisionTree.WellFormed, FABL.CoordinateDecisionTree.wellFormed_queryCoordinates, FABL.DNFFormula, FABL.DNFFormula.canonicalBlockTreeAux, FABL.DNFFormula.canonicalTermIndex, FABL.DNFFormula.compatibleTermIndices, FABL.DNFTerm.Compatible, FABL.DNFTerm.coordinatesAtPositions, FABL.DNFTerm.coordinatesAtPositions_nodup, FABL.DNFTerm.coordinatesAt_freeLiteralPositions_all_free, FABL.DNFTerm.freeLiteralPositions}
\difficulty{5}
Let $n$ be a natural number, let $\varphi$ be a disjunctive normal form on $n$ sign
variables, let $f$ be a natural-number fuel value, and let $\rho$ be a coordinate
restriction. If every coordinate free in $\rho$ belongs to a specified set of available
coordinates, then the tree constructed from $\varphi$, $f$, and $\rho$ is well formed
relative to that set: each queried coordinate is available at its node, and both child
trees are well formed relative to the available set with that coordinate removed.
\end{theorem}

\begin{theorem}[Well-formedness of the raw block tree]
\lean{FABL.DNFFormula.canonicalBlockTreeRaw_wellFormed}
\label{FABL.DNFFormula.canonicalBlockTreeRaw_wellFormed}
\leanok
\uses{FABL.CoordRestriction, FABL.CoordRestriction.freeWithin_univ, FABL.CoordinateDecisionTree.WellFormed, FABL.DNFFormula, FABL.DNFFormula.canonicalBlockTreeAux_wellFormed, FABL.DNFFormula.canonicalBlockTreeRaw}
\difficulty{1}
Let $n$ be a natural number, let $\varphi$ be a disjunctive normal form on $n$ sign
variables, and let $\rho$ be a coordinate restriction. The raw block tree constructed
from $\varphi$ and $\rho$ is well formed relative to the set of all $n$ coordinates:
each queried coordinate belongs to the available set at its node, and both child trees
are well formed relative to the available set with that coordinate removed.
\end{theorem}

\begin{definition}[Canonical block decision tree]
\lean{FABL.DNFFormula.canonicalBlockTree}
\label{FABL.DNFFormula.canonicalBlockTree}
\leanok
\uses{FABL.CoordRestriction, FABL.CoordinateDecisionTree.toF₂DecisionTree, FABL.DNFFormula, FABL.DNFFormula.canonicalBlockTreeRaw, FABL.DNFFormula.canonicalBlockTreeRaw_wellFormed, FABL.DecisionTree, FABL.Sign}
Let $\varphi$ be a disjunctive normal form on $n$ sign variables and let $\rho$ assign
each coordinate a free or fixed sign state. The canonical block decision tree is the raw
canonical block tree for $\varphi$ under $\rho$, represented as a decision tree in which
each query uses an available coordinate and removes it from the coordinates available to
both children.
\end{definition}

\begin{theorem}[Computation by the Canonical Block Tree]
\lean{FABL.DNFFormula.canonicalBlockTree_computes}
\label{FABL.DNFFormula.canonicalBlockTree_computes}
\leanok
\uses{FABL.CoordRestriction, FABL.CoordRestriction.freeCoordinates, FABL.CoordRestriction.restrict, FABL.CoordinateDecisionTree.eval_toF₂DecisionTree, FABL.DNFFormula, FABL.DNFFormula.canonicalBlockTree, FABL.DNFFormula.canonicalBlockTreeRaw, FABL.DNFFormula.eval_canonicalBlockTreeAux, FABL.DNFFormula.toBooleanFunction, FABL.F₂DecisionTree.Computes, FABL.F₂DecisionTree.computes_iff}
\difficulty{3}
Let $n$ be a natural number, let $\varphi$ be a disjunctive normal form on $n$ sign
variables, and let $\rho$ assign to each coordinate one of three states: free, fixed to
$1$, or fixed to $-1$. The decision tree obtained by the canonical block-tree
construction from $\varphi$ and $\rho$ computes the function that evaluates $\varphi$
after completing each free coordinate from a Boolean input and setting each fixed
coordinate to its prescribed sign.
\end{theorem}

\begin{theorem}[Depth bounded by the canonical block tree]
\lean{FABL.DNFFormula.decisionTreeDepth_le_canonicalBlockTree_depth}
\label{FABL.DNFFormula.decisionTreeDepth_le_canonicalBlockTree_depth}
\leanok
\uses{FABL.CoordRestriction, FABL.CoordRestriction.restrict, FABL.DNFFormula, FABL.DNFFormula.canonicalBlockTree, FABL.DNFFormula.canonicalBlockTree_computes, FABL.DNFFormula.toBooleanFunction, FABL.F₂DecisionTree.decisionTreeDepth, FABL.F₂DecisionTree.decisionTreeDepth_le_of_computes, FABL.F₂DecisionTree.depth}
\difficulty{1}
Let $n$ be a nonnegative integer, let $\varphi$ be a DNF formula on $n$ sign
coordinates, and let $\rho$ assign to each coordinate whether it is free or fixed to one
of the two signs. The decision-tree depth of the Boolean function obtained by
restricting $\varphi$ according to $\rho$ is at most the depth of the canonical block
tree of $\varphi$ for $\rho$.
\end{theorem}

\begin{definition}[Switching Block Code]
\lean{FABL.SwitchingBlockCode}
\statementsource{arora-barak-ch13-circuit-lowerbounds}{
  pdf-52929dbaa962-p003-b004,
  pdf-52929dbaa962-p003-b005
}
\label{FABL.SwitchingBlockCode}
\leanok
\uses{FABL.𝔽₂}
For a natural number $w$, a switching block code consists of a set of positions among
$\{0,\ldots,w-1\}$ and an assignment of a bit in $\mathbb{F}_2$ to every position among
them.
\end{definition}

\begin{definition}[Canonical switching block code]
\lean{FABL.SwitchingBlockCode.Canonical}
\label{FABL.SwitchingBlockCode.Canonical}
\leanok
\uses{FABL.SwitchingBlockCode}
A switching block code of width $w$ consists of a set of positions among the $w$
positions and a bit assigned to every position. It is canonical if every position
outside the chosen set is assigned bit $0$.
\end{definition}

\begin{definition}[Nonzero switching block code]
\lean{FABL.SwitchingBlockCode.Nonzero}
\label{FABL.SwitchingBlockCode.Nonzero}
\leanok
\uses{FABL.SwitchingBlockCode}
A switching block code is nonzero if some position in its designated position set has
bit value $1$.
\end{definition}

\begin{definition}[Canonical codes of fixed size]
\lean{FABL.SwitchingBlockCode.FinalOfSize}
\label{FABL.SwitchingBlockCode.FinalOfSize}
\leanok
\uses{FABL.SwitchingBlockCode, FABL.SwitchingBlockCode.Canonical}
For natural numbers $w$ and $i$, a final code of size $i$ is a switching block code on
$w$ positions whose bits are zero outside its position set and whose position set has
cardinality $i$.
\end{definition}

\begin{definition}[Nonzero canonical codes of fixed size]
\lean{FABL.SwitchingBlockCode.FullOfSize}
\label{FABL.SwitchingBlockCode.FullOfSize}
\leanok
\uses{FABL.SwitchingBlockCode, FABL.SwitchingBlockCode.Canonical, FABL.SwitchingBlockCode.Nonzero}
For natural numbers $w$ and $i$, a full code of size $i$ is a switching block code on
$w$ positions whose bits are zero outside its position set, whose position set has
cardinality $i$, and whose bit value is $1$ at some position in that set.
\end{definition}

\begin{definition}[Position sets of fixed size]
\lean{FABL.SwitchingBlockCode.PositionSet}
\label{FABL.SwitchingBlockCode.PositionSet}
\leanok
For natural numbers $w$ and $i$, the position sets of size $i$ are the subsets of
$	ext{Fin}(w)$ having cardinality $i$.
\end{definition}

\begin{theorem}[Count of fixed-size position sets]
\lean{FABL.SwitchingBlockCode.card_positionSet}
\label{FABL.SwitchingBlockCode.card_positionSet}
\leanok
\uses{FABL.SwitchingBlockCode.PositionSet}
\difficulty{1}
For all nonnegative integers $w$ and $i$, the number of $i$-element subsets of a set of
$w$ positions is $\binom{w}{i}$.
\end{theorem}

\begin{definition}[Restriction of code bits]
\lean{FABL.SwitchingBlockCode.restrictBits}
\label{FABL.SwitchingBlockCode.restrictBits}
\leanok
\uses{FABL.SwitchingBlockCode, FABL.𝔽₂}
Given a switching block code $B$, its restricted bit function assigns to each position
in $B$’s position set the bit value that $B$ assigns to that position.
\end{definition}

\begin{definition}[Code from positions and bits]
\lean{FABL.SwitchingBlockCode.ofPositionBits}
\label{FABL.SwitchingBlockCode.ofPositionBits}
\leanok
\uses{FABL.SwitchingBlockCode, FABL.𝔽₂}
Given a finite set $S$ of positions and a bit assignment on $S$, the associated
switching block code has position set $S$, agrees with the assignment on $S$, and has
bit value $0$ at every position outside $S$.
\end{definition}

\begin{theorem}[Positions of a constructed code]
\lean{FABL.SwitchingBlockCode.ofPositionBits_positions}
\label{FABL.SwitchingBlockCode.ofPositionBits_positions}
\leanok
\uses{FABL.SwitchingBlockCode.ofPositionBits, FABL.𝔽₂}
\difficulty{1}
For any set $S$ of positions in a width-$w$ switching block code and any assignment of
bits to the positions in $S$, the code constructed from these data has position set
exactly $S$.
\end{theorem}

\begin{theorem}[Canonical code from position bits]
\lean{FABL.SwitchingBlockCode.canonical_ofPositionBits}
\label{FABL.SwitchingBlockCode.canonical_ofPositionBits}
\leanok
\uses{FABL.SwitchingBlockCode.Canonical, FABL.SwitchingBlockCode.ofPositionBits, FABL.SwitchingBlockCode.ofPositionBits_positions, FABL.𝔽₂}
\difficulty{3}
For any set $S$ of positions in a width-$w$ switching block code and any assignment of
bits to the positions in $S$, the code constructed from these data is canonical: every
position outside $S$ has bit zero.
\end{theorem}

\begin{theorem}[Restriction after position-bit construction]
\lean{FABL.SwitchingBlockCode.restrictBits_ofPositionBits}
\label{FABL.SwitchingBlockCode.restrictBits_ofPositionBits}
\leanok
\uses{FABL.SwitchingBlockCode.ofPositionBits, FABL.SwitchingBlockCode.restrictBits, FABL.𝔽₂}
\difficulty{3}
Let $w$ be a natural number, let $S$ be a finite set of positions in a width-$w$ block,
and assign a bit in $bf_2$ to each position in $S$. The function obtained by
restricting the code constructed from $S$ and these bits to its designated positions
equals the given assignment.
\end{theorem}

\begin{theorem}[Reconstruction of a canonical block code]
\lean{FABL.SwitchingBlockCode.ofPositionBits_restrictBits}
\label{FABL.SwitchingBlockCode.ofPositionBits_restrictBits}
\leanok
\uses{FABL.SwitchingBlockCode, FABL.SwitchingBlockCode.Canonical, FABL.SwitchingBlockCode.ofPositionBits, FABL.SwitchingBlockCode.restrictBits}
\difficulty{4}
Let $B$ be a switching block code whose bits are zero at every position outside its
designated position set. Reconstructing a code from that position set and the
restriction of its bits to it gives back $B$.
\end{theorem}

\begin{definition}[Canonical codes and position bit assignments]
\lean{FABL.SwitchingBlockCode.finalOfSizeEquiv}
\label{FABL.SwitchingBlockCode.finalOfSizeEquiv}
\leanok
\uses{FABL.SwitchingBlockCode.FinalOfSize, FABL.SwitchingBlockCode.PositionSet, FABL.SwitchingBlockCode.canonical_ofPositionBits, FABL.SwitchingBlockCode.ofPositionBits, FABL.SwitchingBlockCode.ofPositionBits_restrictBits, FABL.SwitchingBlockCode.restrictBits, FABL.𝔽₂}
For natural numbers $w$ and $i$, the canonical switching block codes on $w$ positions
whose position sets have cardinality $i$ are in bijection with pairs consisting of an
$i$-element subset of $	ext{Fin}(w)$ and a bit assignment on that subset.
\end{definition}

\begin{theorem}[Count of canonical codes of fixed size]
\lean{FABL.SwitchingBlockCode.card_finalOfSize}
\label{FABL.SwitchingBlockCode.card_finalOfSize}
\leanok
\uses{FABL.SwitchingBlockCode.FinalOfSize, FABL.SwitchingBlockCode.PositionSet, FABL.SwitchingBlockCode.card_positionSet, FABL.SwitchingBlockCode.finalOfSizeEquiv, FABL.𝔽₂}
\difficulty{4}
For all nonnegative integers $w$ and $i$, the number of canonical width-$w$ switching
block codes whose position set has size $i$ is $\binom{w}{i}2^i$.
\end{theorem}

\begin{theorem}[Count of nonzero bit assignments]
\lean{FABL.SwitchingBlockCode.card_nonzero_positionBits}
\label{FABL.SwitchingBlockCode.card_nonzero_positionBits}
\leanok
\uses{FABL.𝔽₂}
\difficulty{3}
For any finite set $S$ of positions, the number of nonzero functions from $S$ to the
two-element field is $2^{|S|}-1$.
\end{theorem}

\begin{theorem}[Nonzero code and restricted bits]
\lean{FABL.SwitchingBlockCode.nonzero_iff_restrictBits_ne_zero}
\label{FABL.SwitchingBlockCode.nonzero_iff_restrictBits_ne_zero}
\leanok
\uses{FABL.SwitchingBlockCode, FABL.SwitchingBlockCode.Nonzero, FABL.SwitchingBlockCode.restrictBits}
\difficulty{5}
For any width-$w$ switching block code, there is a position in its position set with bit
$1$ if and only if the restriction of its bit assignment to that position set is not the
zero function.
\end{theorem}

\begin{definition}[Nonzero codes and position bit assignments]
\lean{FABL.SwitchingBlockCode.fullOfSizeEquiv}
\label{FABL.SwitchingBlockCode.fullOfSizeEquiv}
\leanok
\uses{FABL.SwitchingBlockCode.FullOfSize, FABL.SwitchingBlockCode.PositionSet, FABL.SwitchingBlockCode.canonical_ofPositionBits, FABL.SwitchingBlockCode.nonzero_iff_restrictBits_ne_zero, FABL.SwitchingBlockCode.ofPositionBits, FABL.SwitchingBlockCode.ofPositionBits_restrictBits, FABL.SwitchingBlockCode.restrictBits, FABL.SwitchingBlockCode.restrictBits_ofPositionBits, FABL.𝔽₂}
For natural numbers $w$ and $i$, the canonical switching block codes on $w$ positions
whose position sets have cardinality $i$ and whose bits are not all zero on that set are
in bijection with pairs consisting of an $i$-element subset of $	ext{Fin}(w)$ and a
nonzero bit assignment on that subset.
\end{definition}

\begin{theorem}[Count of nonzero canonical codes]
\lean{FABL.SwitchingBlockCode.card_fullOfSize}
\label{FABL.SwitchingBlockCode.card_fullOfSize}
\leanok
\uses{FABL.SwitchingBlockCode.FullOfSize, FABL.SwitchingBlockCode.PositionSet, FABL.SwitchingBlockCode.card_nonzero_positionBits, FABL.SwitchingBlockCode.card_positionSet, FABL.SwitchingBlockCode.fullOfSizeEquiv, FABL.𝔽₂}
\difficulty{4}
For all nonnegative integers $w$ and $i$, the number of canonical width-$w$ switching
block codes whose position set has size $i$ and whose bits are not all zero on that set
is $\binom{w}{i}(2^i-1)$.
\end{theorem}

\begin{definition}[Switching codes]
\lean{FABL.SwitchingCode}
\label{FABL.SwitchingCode}
\leanok
\uses{FABL.SwitchingBlockCode.FinalOfSize, FABL.SwitchingBlockCode.FullOfSize}
For a width $w$ and a natural number $m$, a switching code of index $m$ is defined
inductively as follows. At index $k+1$, it is either a canonical block code whose
position set has size $k+1$, or a choice of $i<k$, a canonical nonzero block code whose
position set has size $i+1$, and a switching code of index $k-i$.
\end{definition}

\begin{definition}[Zero-index switching codes]
\lean{FABL.SwitchingCode.zeroEquivEmpty}
\label{FABL.SwitchingCode.zeroEquivEmpty}
\leanok
\uses{FABL.SwitchingCode}
For every width $w$, the type of switching codes of index $0$ is equivalent to the empty
type; in particular, there are no such codes.
\end{definition}

\begin{definition}[Successor decomposition of switching codes]
\lean{FABL.SwitchingCode.succEquiv}
\label{FABL.SwitchingCode.succEquiv}
\leanok
\uses{FABL.SwitchingBlockCode.FinalOfSize, FABL.SwitchingBlockCode.FullOfSize, FABL.SwitchingCode}
For every width $w$ and natural number $k$, switching codes of index $k+1$ are
equivalent to the disjoint union of canonical block codes with $k+1$ positions and
tuples consisting of an integer $i<k$, a canonical nonzero block code with $i+1$
positions, and a switching code of index $k-i$.
\end{definition}

\begin{definition}[Finite switching-code types]
\lean{FABL.SwitchingCode.fintype}
\label{FABL.SwitchingCode.fintype}
\leanok
\uses{FABL.SwitchingBlockCode.FinalOfSize, FABL.SwitchingBlockCode.FullOfSize, FABL.SwitchingCode, FABL.SwitchingCode.succEquiv, FABL.SwitchingCode.zeroEquivEmpty}
For every width $w$, each type of switching codes of index $k$, for any natural number
$k$, is finite.
\end{definition}

\begin{theorem}[Cardinality recurrence for switching codes]
\lean{FABL.SwitchingCode.card_succ}
\label{FABL.SwitchingCode.card_succ}
\leanok
\uses{FABL.SwitchingBlockCode.card_finalOfSize, FABL.SwitchingBlockCode.card_fullOfSize, FABL.SwitchingCode, FABL.SwitchingCode.succEquiv}
\difficulty{3}
For all nonnegative integers $w$ and $k$, the number of switching codes of width $w$ and
index $k+1$ equals
$\binom{w}{k+1}2^{k+1}+\sum_{i=0}^{k-1}\binom{w}{i+1}(2^{i+1}-1)\,C(w,k-i)$, where
$C(w,j)$ denotes the number of switching codes of width $w$ and index $j$.
\end{theorem}

\begin{definition}[Full-block coefficient]
\lean{FABL.SwitchingCode.fullCoefficient}
\label{FABL.SwitchingCode.fullCoefficient}
\leanok
For each natural number $i$, define the real coefficient $a_i=(2^i-1)(4/9)^i/i!$.
\end{definition}

\begin{definition}[Final-block remainder]
\lean{FABL.SwitchingCode.finalRemainder}
\label{FABL.SwitchingCode.finalRemainder}
\leanok
For each natural number $i$, define the real remainder $r_i=(4/9)^i/i!$.
\end{definition}

\begin{theorem}[Factorial lower bound]
\lean{FABL.SwitchingCode.factorial_ge_twentyFour_mul_four_pow}
\label{FABL.SwitchingCode.factorial_ge_twentyFour_mul_four_pow}
\leanok
\difficulty{4}
For every integer $i \ge 4$, $24\cdot 4^{i-4} \le i!$.
\end{theorem}

\begin{theorem}[Full coefficient tail bound]
\lean{FABL.SwitchingCode.fullCoefficient_tail}
\label{FABL.SwitchingCode.fullCoefficient_tail}
\leanok
\uses{FABL.SwitchingCode.factorial_ge_twentyFour_mul_four_pow, FABL.SwitchingCode.fullCoefficient}
\difficulty{5}
For every integer $i \ge 4$, let $c_i=((2^i-1)(4/9)^i)/i!$. Then $c_i \le
(32/3)(2/9)^i$.
\end{theorem}

\begin{theorem}[Geometric sum bound]
\lean{FABL.SwitchingCode.geometric_two_ninths_bound}
\label{FABL.SwitchingCode.geometric_two_ninths_bound}
\leanok
\difficulty{4}
For every integer $m \ge 0$, $\sum_{j=0}^{m-1}(2/9)^j \le 9/7$, where the sum is empty
when $m=0$.
\end{theorem}

\begin{theorem}[Bound on the coefficient sum]
\lean{FABL.SwitchingCode.sum_fullCoefficient_range}
\label{FABL.SwitchingCode.sum_fullCoefficient_range}
\leanok
\uses{FABL.SwitchingCode.fullCoefficient, FABL.SwitchingCode.fullCoefficient_tail, FABL.SwitchingCode.geometric_two_ninths_bound}
\difficulty{5}
For every integer $k\ge 0$, define $c_i=((2^i-1)(4/9)^i)/i!$. Then
$\sum_{i=0}^{k-1}c_{i+1}\le 13420/15309$, with the sum empty when $k=0$.
\end{theorem}

\begin{theorem}[Final remainder bound]
\lean{FABL.SwitchingCode.finalRemainder_le}
\label{FABL.SwitchingCode.finalRemainder_le}
\leanok
\uses{FABL.SwitchingCode.finalRemainder}
\difficulty{4}
For every integer $i \ge 2$, the quantity $((4/9)^i)/i!$ is at most $8/81$.
\end{theorem}

\begin{theorem}[Scaled full coefficient identity]
\lean{FABL.SwitchingCode.scale_mul_fullCoefficient}
\label{FABL.SwitchingCode.scale_mul_fullCoefficient}
\leanok
\uses{FABL.SwitchingCode.fullCoefficient}
\difficulty{4}
For all natural numbers $w$ and $i$, define $c_i=((2^i-1)(4/9)^i)/i!$. Then
$((9w/4)^i)c_i=(w^i/i!)(2^i-1)$.
\end{theorem}

\begin{theorem}[Scaled coefficient identity]
\lean{FABL.SwitchingCode.scale_mul_finalCoefficient}
\label{FABL.SwitchingCode.scale_mul_finalCoefficient}
\leanok
\uses{FABL.SwitchingCode.finalRemainder, FABL.SwitchingCode.fullCoefficient}
\difficulty{4}
For all natural numbers $w$ and $i$, define $c_i=((2^i-1)(4/9)^i)/i!$ and
$r_i=(4/9)^i/i!$. Then $((9w/4)^i)(c_i+r_i)=(w^i/i!)2^i$.
\end{theorem}

\begin{theorem}[Bound for the full-block term]
\lean{FABL.SwitchingCode.choose_full_le}
\label{FABL.SwitchingCode.choose_full_le}
\leanok
\uses{FABL.SwitchingCode.fullCoefficient, FABL.SwitchingCode.scale_mul_fullCoefficient}
\difficulty{4}
For all nonnegative integers $w$ and $i$, define $a_i=((2^i-1)(4/9)^i)/i!$. Then
$\binom{w}{i}(2^i-1)\leq (9w/4)^i a_i$.
\end{theorem}

\begin{theorem}[Bound for the final-block term]
\lean{FABL.SwitchingCode.choose_final_le}
\label{FABL.SwitchingCode.choose_final_le}
\leanok
\uses{FABL.SwitchingCode.finalRemainder, FABL.SwitchingCode.fullCoefficient, FABL.SwitchingCode.scale_mul_finalCoefficient}
\difficulty{4}
For all nonnegative integers $w$ and $i$, define $a_i=((2^i-1)(4/9)^i)/i!$ and
$r_i=(4/9)^i/i!$. Then $\binom{w}{i}2^i\leq (9w/4)^i(a_i+r_i)$.
\end{theorem}

\begin{theorem}[Switching-code cardinality bound]
\lean{FABL.SwitchingCode.card_le_scale_pow}
\label{FABL.SwitchingCode.card_le_scale_pow}
\leanok
\uses{FABL.SwitchingCode, FABL.SwitchingCode.card_succ, FABL.SwitchingCode.choose_final_le, FABL.SwitchingCode.choose_full_le, FABL.SwitchingCode.finalRemainder, FABL.SwitchingCode.finalRemainder_le, FABL.SwitchingCode.fullCoefficient, FABL.SwitchingCode.sum_fullCoefficient_range, FABL.SwitchingCode.zeroEquivEmpty}
\difficulty{6}
For all nonnegative integers $w$ and $k$, the number of switching codes of width $w$ and
index $k$ is at most $\left(\frac{9w}{4}\right)^k$.
\end{theorem}

\begin{definition}[Extension of coordinate restrictions]
\lean{FABL.CoordRestriction.Extends}
\label{FABL.CoordRestriction.Extends}
\leanok
\uses{FABL.CoordRestriction}
Let $n$ be a natural number, and let $
ho,eta$ be coordinate restrictions on
${0,1,-1}^n$, where each coordinate is free, fixed to $1$, or fixed to $-1$. Say that
$eta$ extends $rho$ if every coordinate that $rho$ fixes has the same value in
$eta$.
\end{definition}

\begin{theorem}[Reflexivity of extension]
\lean{FABL.CoordRestriction.extends_refl}
\label{FABL.CoordRestriction.extends_refl}
\leanok
\uses{FABL.CoordRestriction, FABL.CoordRestriction.Extends}
\difficulty{1}
Let $n$ be a natural number and let $\rho$ assign a restriction state to each of the $n$
coordinates. Then $\rho$ agrees with itself at every coordinate where it is not free.
\end{theorem}

\begin{theorem}[Transitivity of restriction extension]
\lean{FABL.CoordRestriction.Extends.trans}
\label{FABL.CoordRestriction.Extends.trans}
\leanok
\uses{FABL.CoordRestriction, FABL.CoordRestriction.Extends}
\difficulty{4}
Let $n$ be a natural number and let $\rho,\eta,\tau$ be coordinate restrictions on $n$
coordinates. If every coordinate fixed by $\rho$ has the same state in $\eta$, and every
coordinate fixed by $\eta$ has the same state in $\tau$, then every coordinate fixed by
$\rho$ has the same state in $\tau$.
\end{theorem}

\begin{theorem}[Agreement on a fixed coordinate]
\lean{FABL.CoordRestriction.Extends.eq_of_ne_free}
\label{FABL.CoordRestriction.Extends.eq_of_ne_free}
\leanok
\uses{FABL.CoordRestriction, FABL.CoordRestriction.Extends}
\difficulty{1}
Let $\rho,\eta:\{1,\ldots,n\}\to\{\mathrm{free},\mathrm{fix\ one},\mathrm{fix\ negative\
one}\}$, where $n$ is a natural number. If $\eta$ extends $\rho$ in the sense that it
agrees with $\rho$ at every coordinate where $\rho$ is not free, then $\eta(i)=\rho(i)$
for every coordinate $i$ at which $\rho(i)$ is not free.
\end{theorem}

\begin{theorem}[Extension after fixing a list of coordinates]
\lean{FABL.CoordRestriction.extends_fixCoordinateList}
\label{FABL.CoordRestriction.extends_fixCoordinateList}
\leanok
\uses{FABL.CoordRestriction, FABL.CoordRestriction.Extends, FABL.CoordRestriction.fixCoordinateList, FABL.CoordRestriction.fixCoordinateList_apply_of_not_mem}
\difficulty{4}
Let $n$ be a natural number, let $\rho$ assign a restriction state to each of the $n$
coordinates, let $L$ be a list of coordinates, and let $x\in\mathbb{F}_2^n$. If every
coordinate in $L$ is free in $\rho$, then fixing the coordinates in $L$ according to the
signs encoded by $x$ produces an assignment that agrees with $\rho$ at every coordinate
where $\rho$ was not free.
\end{theorem}

\begin{theorem}[Compatibility under restriction extension]
\lean{FABL.DNFTerm.compatible_of_extends}
\label{FABL.DNFTerm.compatible_of_extends}
\leanok
\uses{FABL.CoordRestriction, FABL.CoordRestriction.Extends, FABL.CoordRestriction.complete, FABL.DNFTerm, FABL.DNFTerm.Compatible, FABL.binarySignEquiv}
\difficulty{5}
Let $n$ be a natural number, let $T$ be a DNF term on $n$ sign coordinates, and let
$\rho,\eta$ be coordinate restrictions such that every coordinate fixed by $\rho$ has
the same fixed value in $\eta$. If $T$ is compatible with $\eta$, then $T$ is compatible
with $\rho$.
\end{theorem}

\begin{theorem}[Compatibility under extension]
\lean{FABL.DNFTerm.compatible_of_extends_of_no_free}
\label{FABL.DNFTerm.compatible_of_extends_of_no_free}
\leanok
\uses{FABL.CoordRestriction, FABL.CoordRestriction.Extends, FABL.DNFTerm, FABL.DNFTerm.Compatible, FABL.DNFTerm.eval_eq_neg_one_iff, FABL.DNFTerm.freeLiteralPositions, FABL.DNFTerm.mem_freeLiteralPositions_iff, FABL.signEncode}
\difficulty{5}
Let $n$ be a natural number, let $T$ be a DNF term on $n$ sign coordinates, and let
$\rho,\eta$ be coordinate restrictions. Suppose $T$ is compatible with $\rho$, every
coordinate appearing in a literal of $T$ is fixed by $\rho$, and $\eta$ agrees with
$\rho$ at every coordinate where $\rho$ is not free. Then $T$ is compatible with $\eta$.
\end{theorem}

\begin{theorem}[Canonical index under an extending restriction]
\lean{FABL.DNFFormula.canonicalTermIndex_eq_of_extends}
\label{FABL.DNFFormula.canonicalTermIndex_eq_of_extends}
\leanok
\uses{FABL.CoordRestriction, FABL.CoordRestriction.Extends, FABL.DNFFormula, FABL.DNFFormula.canonicalTermIndex, FABL.DNFFormula.compatibleTermIndices, FABL.DNFTerm.Compatible, FABL.DNFTerm.compatible_of_extends}
\difficulty{4}
Let $\phi$ be a DNF formula on $n$ coordinates, and let $\rho,\eta$ be coordinate
restrictions such that $\eta$ agrees with $\rho$ wherever $\rho$ is fixed. Suppose each
restriction admits at least one compatible term, and the term of least index compatible
with $\rho$ is also compatible with $\eta$. Then the least index of a term compatible
with $\eta$ equals the least index of a term compatible with $\rho$.
\end{theorem}

\begin{definition}[Free Coordinates in a List]
\lean{FABL.CoordRestriction.freeCoordinateList}
\statementsource{Analysis-of-Boolean-Functions-by-Ryan-ODonnell}{
  pdf-ffbe9c7b3306-p098-b003,
  pdf-ffbe9c7b3306-p098-b004
}
\label{FABL.CoordRestriction.freeCoordinateList}
\leanok
\uses{FABL.CoordRestriction}
Let $\eta$ assign each of the $n$ coordinates a state that is either free, fixed at $1$,
or fixed at $-1$, and let $L$ be a list of coordinates. The resulting state assignment
makes coordinate $i$ free if $i$ occurs in $L$, and otherwise leaves its state as
$\eta(i)$.
\end{definition}

\begin{theorem}[Listed coordinates are free]
\lean{FABL.CoordRestriction.freeCoordinateList_apply_of_mem}
\label{FABL.CoordRestriction.freeCoordinateList_apply_of_mem}
\leanok
\uses{FABL.CoordRestriction, FABL.CoordRestriction.freeCoordinateList}
\difficulty{1}
Let $n$ be a natural number, let $\eta$ assign one of the three restriction states to
each of the $n$ coordinates, and let $L$ be a list of coordinates. Every coordinate
belonging to $L$ is free in the assignment obtained by making the coordinates in $L$
free.
\end{theorem}

\begin{theorem}[Free coordinate list outside the list]
\lean{FABL.CoordRestriction.freeCoordinateList_apply_of_not_mem}
\label{FABL.CoordRestriction.freeCoordinateList_apply_of_not_mem}
\leanok
\uses{FABL.CoordRestriction, FABL.CoordRestriction.freeCoordinateList}
\difficulty{1}
Let $n$ be a natural number, let $\eta$ assign one of the states free, fixed at $1$, or
fixed at $-1$ to each of the $n$ coordinates, and let $L$ be a list of coordinates. For
any coordinate $i$ not in $L$, the restriction obtained by making every coordinate in
$L$ free agrees with $\eta$ at $i$.
\end{theorem}

\begin{theorem}[Free coordinates after fixing a list]
\lean{FABL.CoordRestriction.freeCoordinateList_fixCoordinateList}
\label{FABL.CoordRestriction.freeCoordinateList_fixCoordinateList}
\leanok
\uses{FABL.CoordRestriction, FABL.CoordRestriction.fixCoordinateList, FABL.CoordRestriction.fixCoordinateList_apply_of_not_mem, FABL.CoordRestriction.freeCoordinateList, FABL.CoordRestriction.freeCoordinateList_apply_of_mem, FABL.CoordRestriction.freeCoordinateList_apply_of_not_mem}
\difficulty{4}
Let $n$ be a natural number, let $\rho$ assign a restriction state to each of the $n$
coordinates, let $L$ be a list of coordinates, and let $x$ assign a bit to each
coordinate. If every coordinate in $L$ is free under $\rho$, then fixing the coordinates
in $L$ according to $x$ and making all coordinates in $L$ free again yields exactly
$\rho$.
\end{theorem}

\begin{theorem}[Fixing an already matched list]
\lean{FABL.CoordRestriction.fixCoordinateList_eq_self}
\label{FABL.CoordRestriction.fixCoordinateList_eq_self}
\leanok
\uses{FABL.CoordRestriction, FABL.CoordRestriction.fixCoordinate, FABL.CoordRestriction.fixCoordinateList, FABL.CoordRestriction.fixCoordinateList_cons, FABL.CoordRestriction.fixCoordinate_apply_of_ne, FABL.CoordRestriction.fixedState, FABL.signEncode}
\difficulty{5}
Let $n$ be a natural number, let $\eta$ assign a restriction state to each of the $n$
coordinates, let $L$ be a list of coordinates, and let $x\in\mathbb{F}_2^n$. If at every
coordinate in $L$, $\eta$ already equals the fixed state determined by the sign encoded
by $x_i$, then fixing the coordinates in $L$ according to $x$ leaves $\eta$ unchanged.
\end{theorem}

\begin{theorem}[Value at a listed coordinate]
\lean{FABL.CoordRestriction.fixCoordinateList_apply_of_mem}
\label{FABL.CoordRestriction.fixCoordinateList_apply_of_mem}
\leanok
\uses{FABL.CoordRestriction, FABL.CoordRestriction.fixCoordinateList, FABL.CoordRestriction.fixCoordinateList_apply_of_not_mem, FABL.CoordRestriction.fixCoordinateList_cons, FABL.CoordRestriction.fixedState, FABL.signEncode}
\difficulty{4}
Let $n$ be a natural number, let $\eta$ assign a restriction state to each of the $n$
coordinates, let $L$ be a list of coordinates with no repetitions, and let
$x\in\mathbb{F}_2^n$. For any coordinate $i$ in $L$, the assignment obtained by fixing
the coordinates in $L$ according to the signs encoded by $x$ has at $i$ the fixed state
determined by the sign encoded by $x_i$.
\end{theorem}

\begin{theorem}[Extension after parallel coordinate fixing]
\lean{FABL.CoordRestriction.Extends.parallel_fixCoordinateList}
\label{FABL.CoordRestriction.Extends.parallel_fixCoordinateList}
\leanok
\uses{FABL.CoordRestriction, FABL.CoordRestriction.Extends, FABL.CoordRestriction.fixCoordinateList, FABL.CoordRestriction.fixCoordinateList_apply_of_mem, FABL.CoordRestriction.fixCoordinateList_apply_of_not_mem}
\difficulty{4}
Let $n$ be a natural number, let
$\rho,\tau:\{1,\ldots,n\}\to\{\mathrm{free},\mathrm{fix\ one},\mathrm{fix\ negative\
one}\}$, and let $L$ be a list of distinct coordinates in $\{1,\ldots,n\}$. For
$x,y\in(\mathbb F_2)^n$, suppose fixing the coordinates in $L$ in $\rho$ according to
the signs encoded by $x$ gives a restriction extended by $\tau$. Then fixing those
coordinates in both $\rho$ and $\tau$ according to the signs encoded by $y$ gives
restrictions such that the one obtained from $\tau$ extends the one obtained from
$\rho$.
\end{theorem}

\begin{theorem}[Repeated coordinate updates are overwritten]
\lean{FABL.CoordRestriction.fixCoordinateList_overwrite}
\label{FABL.CoordRestriction.fixCoordinateList_overwrite}
\leanok
\uses{FABL.CoordRestriction, FABL.CoordRestriction.fixCoordinateList, FABL.CoordRestriction.fixCoordinateList_apply_of_mem, FABL.CoordRestriction.fixCoordinateList_apply_of_not_mem}
\difficulty{4}
Let $n$ be a natural number, let $\eta$ assign one of the three restriction states to
each of the $n$ coordinates, let $L$ be a list of distinct coordinates, and let
$x,y\in(\mathbb{Z}/2\mathbb{Z})^n$. Applying the updates prescribed by $x$ on $L$,
followed by those prescribed by $y$ on $L$, has the same result as applying only the
updates prescribed by $y$ on $L$.
\end{theorem}

\begin{definition}[Position embedding for a DNF term]
\lean{FABL.DNFTerm.positionEmbedding}
\label{FABL.DNFTerm.positionEmbedding}
\leanok
\uses{FABL.DNFTerm, FABL.DNFTerm.width}
Let $T$ be a DNF term on $n$ coordinates, with distinct literal indices, and let $w$ be
a natural number such that the width of $T$ is at most $w$. Then the literal positions
of $T$ embed injectively into the first $w$ positions, preserving their indices.
\end{definition}

\begin{definition}[Position of a Supported Coordinate]
\lean{FABL.DNFTerm.positionOfSupport}
\label{FABL.DNFTerm.positionOfSupport}
\leanok
\uses{FABL.DNFTerm, FABL.DNFTerm.literalAt_mem, FABL.DNFTerm.support}
Let $T$ be a DNF term on $n$ variables, and let $i$ be a coordinate in its support. The
position of $i$ is the index in the term’s literal list of the literal whose variable is
$i$.
\end{definition}

\begin{theorem}[Literal at a support coordinate]
\lean{FABL.DNFTerm.get_positionOfSupport}
\label{FABL.DNFTerm.get_positionOfSupport}
\leanok
\uses{FABL.DNFTerm, FABL.DNFTerm.literalAt, FABL.DNFTerm.literalAt_mem, FABL.DNFTerm.positionOfSupport, FABL.DNFTerm.support}
\difficulty{1}
Let $T$ be a conjunction of literals on $n$ coordinates with pairwise distinct literal
coordinates, and let $i$ be a coordinate in the support of $T$. The literal selected
from $T$ at coordinate $i$ is the literal at the corresponding position in the list of
literals.
\end{theorem}

\begin{theorem}[Coordinate of the selected literal]
\lean{FABL.DNFTerm.positionOfSupport_index}
\label{FABL.DNFTerm.positionOfSupport_index}
\leanok
\uses{FABL.DNFTerm, FABL.DNFTerm.get_positionOfSupport, FABL.DNFTerm.literalAt_index, FABL.DNFTerm.positionOfSupport, FABL.DNFTerm.support}
\difficulty{1}
Let $n$ be a natural number and let $T$ be a DNF term on $n$ coordinates. For any
coordinate $i$ in the support of $T$, the literal selected from $T$ at coordinate $i$
has coordinate $i$.
\end{theorem}

\begin{theorem}[Position of a supported literal]
\lean{FABL.DNFTerm.positionOfSupport_literal}
\label{FABL.DNFTerm.positionOfSupport_literal}
\leanok
\uses{FABL.DNFTerm, FABL.DNFTerm.mem_support_of_mem_literals, FABL.DNFTerm.positionOfSupport, FABL.DNFTerm.positionOfSupport_index}
\difficulty{4}
Let $n$ be a natural number and let $T$ be a DNF term on $n$ coordinates. For every
valid position $j$ in its literal list, the position selected for the coordinate of the
literal at $j$ is exactly $j$.
\end{theorem}

\begin{definition}[Term Positions in a Block Code]
\lean{FABL.SwitchingBlockCode.termPositions}
\label{FABL.SwitchingBlockCode.termPositions}
\leanok
\uses{FABL.DNFTerm, FABL.DNFTerm.positionEmbedding, FABL.DNFTerm.width, FABL.SwitchingBlockCode}
Let $B$ be a switching block code of length $w$, let $T$ be a DNF term on $n$ variables,
and suppose the width of $T$ is at most $w$. The term positions are the positions in $T$
whose images among the first $w$ positions belong to the position set of $B$.
\end{definition}

\begin{theorem}[Positions of a mapped term code]
\lean{FABL.SwitchingBlockCode.termPositions_of_mapped}
\label{FABL.SwitchingBlockCode.termPositions_of_mapped}
\leanok
\uses{FABL.DNFTerm, FABL.DNFTerm.positionEmbedding, FABL.DNFTerm.width, FABL.SwitchingBlockCode.ofPositionBits, FABL.SwitchingBlockCode.termPositions, FABL.𝔽₂}
\difficulty{1}
Let $T$ be a DNF term of width at most $w$, and let $S$ be a set of positions in the
list of literals of $T$. Map these positions into the first $w$ coordinates using their
natural coordinate embedding, and construct a switching block code with that mapped
position set and any bit assignment on it. The positions of this code within $T$ are
exactly $S$.
\end{theorem}

\begin{definition}[Relative Bit of a Literal]
\lean{FABL.DNFTerm.relativeBit}
\label{FABL.DNFTerm.relativeBit}
\leanok
\uses{FABL.DNFTerm, FABL.binarySignEquiv, FABL.𝔽₂}
Let $T$ be a DNF term on $n$ variables, let $x$ be a binary assignment to those
variables, and let $j$ be a position in the literal list of $T$. The relative bit at $j$
is the sum modulo $2$ of the bit of $x$ at that literal’s variable and the bit encoding
the sign required by the literal.
\end{definition}

\begin{definition}[Block Code of a DNF Term]
\lean{FABL.DNFTerm.blockCode}
\label{FABL.DNFTerm.blockCode}
\leanok
\uses{FABL.DNFTerm, FABL.DNFTerm.positionEmbedding, FABL.DNFTerm.relativeBit, FABL.DNFTerm.width, FABL.SwitchingBlockCode, FABL.SwitchingBlockCode.ofPositionBits}
Let $T$ be a DNF term on $n$ variables whose width is at most $w$, let $P$ be a set of
positions in $T$, and let $x$ be a binary assignment to the variables. The block code
has as its positions the images of the elements of $P$ among the $w$ positions, and at
each such position its bit is the relative bit of $x$ at the corresponding literal.
\end{definition}

\begin{theorem}[Block code position set]
\lean{FABL.DNFTerm.blockCode_positions}
\label{FABL.DNFTerm.blockCode_positions}
\leanok
\uses{FABL.DNFTerm, FABL.DNFTerm.blockCode, FABL.DNFTerm.positionEmbedding, FABL.DNFTerm.width}
\difficulty{1}
Let $T$ be a conjunction of literals on $n$ Boolean coordinates with width at most $w$,
and let $x\in\mathbb{F}_2^n$. The positions of the block code formed from a set $S$ of
literal positions are exactly the images of $S$ under the natural embedding of the
term's positions into the first $w$ positions.
\end{theorem}

\begin{theorem}[Canonical block code]
\lean{FABL.DNFTerm.blockCode_canonical}
\label{FABL.DNFTerm.blockCode_canonical}
\leanok
\uses{FABL.DNFTerm, FABL.DNFTerm.blockCode, FABL.DNFTerm.width, FABL.SwitchingBlockCode.Canonical, FABL.SwitchingBlockCode.canonical_ofPositionBits}
\difficulty{1}
Let $T$ be a conjunction of literals on distinct coordinates with width at most $w$, let
$P$ be any set of its literal positions, and let $x\in\mathbb{F}_2^n$. The block code
formed from these data is canonical: every bit at a position outside its selected
positions is zero.
\end{theorem}

\begin{theorem}[Block code recovers term positions]
\lean{FABL.DNFTerm.blockCode_termPositions}
\label{FABL.DNFTerm.blockCode_termPositions}
\leanok
\uses{FABL.DNFTerm, FABL.DNFTerm.blockCode, FABL.DNFTerm.width, FABL.SwitchingBlockCode.termPositions, FABL.SwitchingBlockCode.termPositions_of_mapped}
\difficulty{1}
Let $T$ be a conjunction of literals on $n$ Boolean coordinates with width at most $w$,
and let $x\in\mathbb{F}_2^n$. Regard the block code formed from a set $S$ of literal
positions as a set of positions among the first $w$ coordinates. The literal positions
corresponding to those code positions are exactly $S$.
\end{theorem}

\begin{theorem}[Block code position count]
\lean{FABL.DNFTerm.blockCode_card_positions}
\label{FABL.DNFTerm.blockCode_card_positions}
\leanok
\uses{FABL.DNFTerm, FABL.DNFTerm.blockCode, FABL.DNFTerm.blockCode_positions, FABL.DNFTerm.width}
\difficulty{1}
Let $n,w\in\mathbb{N}$, let $T$ be a conjunction of literals on $n$ Boolean coordinates,
and let $x\in\mathbb{F}_2^n$. If the number of literals in $T$ is at most $w$, then the
block code formed from any set of literal positions has exactly as many positions as
that set.
\end{theorem}

\begin{definition}[Assignment from a Block Code]
\lean{FABL.DNFTerm.assignmentOfBlock}
\label{FABL.DNFTerm.assignmentOfBlock}
\leanok
\uses{FABL.DNFTerm, FABL.DNFTerm.positionEmbedding, FABL.DNFTerm.positionOfSupport, FABL.DNFTerm.requiredAt, FABL.DNFTerm.support, FABL.DNFTerm.width, FABL.SwitchingBlockCode, FABL.binarySignEquiv}
Let $T$ be a DNF term on $n$ variables whose width is at most $w$, and let $B$ be a
switching block code of length $w$. The resulting binary assignment is zero outside the
support of $T$; on each supported coordinate, it is the sum modulo $2$ of the bit of $B$
at the image of that coordinate’s position in $T$ and the bit encoding the sign required
by its literal.
\end{definition}

\begin{theorem}[Assignment at a literal coordinate]
\lean{FABL.DNFTerm.assignmentOfBlock_literal}
\label{FABL.DNFTerm.assignmentOfBlock_literal}
\leanok
\uses{FABL.DNFTerm, FABL.DNFTerm.assignmentOfBlock, FABL.DNFTerm.literalAt_eq, FABL.DNFTerm.mem_support_of_mem_literals, FABL.DNFTerm.positionEmbedding, FABL.DNFTerm.positionOfSupport_literal, FABL.DNFTerm.requiredAt, FABL.DNFTerm.width, FABL.SwitchingBlockCode, FABL.binarySignEquiv}
\difficulty{4}
Let $T$ be a conjunction of literals on distinct coordinates with width at most $w$, and
let $B$ be a block code on $w$ positions. At the coordinate of any literal in $T$, the
assignment constructed from $B$ equals the bit of $B$ at that literal's embedded
position, added in $\mathbb{F}_2$ to the inverse sign encoding of the literal's required
sign.
\end{theorem}

\begin{theorem}[Bits of the block code]
\lean{FABL.DNFTerm.blockCode_bits}
\label{FABL.DNFTerm.blockCode_bits}
\leanok
\uses{FABL.DNFTerm, FABL.DNFTerm.blockCode, FABL.DNFTerm.positionEmbedding, FABL.DNFTerm.relativeBit, FABL.DNFTerm.width, FABL.SwitchingBlockCode.ofPositionBits}
\difficulty{4}
Let $T$ be a conjunction of literals on distinct coordinates with width at most $w$, let
$P$ be a set of its literal positions, and let $x\in\mathbb{F}_2^n$. The block code
formed from $T$, $P$, and $x$ has, at the embedded position of every $j\in P$, the bit
of $x$ at the coordinate of the $j$th literal, adjusted by the inverse sign encoding of
that literal's required sign.
\end{theorem}

\begin{theorem}[Assignment from its block code]
\lean{FABL.DNFTerm.assignmentOfBlock_blockCode_literal}
\label{FABL.DNFTerm.assignmentOfBlock_blockCode_literal}
\leanok
\uses{FABL.DNFTerm, FABL.DNFTerm.assignmentOfBlock, FABL.DNFTerm.assignmentOfBlock_literal, FABL.DNFTerm.blockCode, FABL.DNFTerm.blockCode_bits, FABL.DNFTerm.relativeBit, FABL.DNFTerm.width}
\difficulty{3}
Let $T$ be a conjunction of literals on distinct coordinates with width at most $w$. For
any block code $B$ on $w$ positions and any literal position $j$ in $T$, the assignment
constructed from $B$ for $T$ at that literal's coordinate equals the bit of $B$ at the
embedded position of $j$, adjusted by the inverse sign encoding of that literal's
required sign.
\end{theorem}

\begin{definition}[Coordinates Selected by a Block Code]
\lean{FABL.SwitchingBlockCode.termCoordinates}
\label{FABL.SwitchingBlockCode.termCoordinates}
\leanok
\uses{FABL.DNFTerm, FABL.DNFTerm.coordinatesAtPositions, FABL.SwitchingBlockCode}
Let $B$ be a switching block code of length $w$ and let $T$ be a DNF term on $n$
variables. The term coordinates are the variable coordinates of those literals whose
positions, when viewed among the first $w$ positions, belong to the position set of $B$.
\end{definition}

\begin{theorem}[Term coordinates of its block code]
\lean{FABL.DNFTerm.termCoordinates_blockCode}
\label{FABL.DNFTerm.termCoordinates_blockCode}
\leanok
\uses{FABL.DNFTerm, FABL.DNFTerm.blockCode, FABL.DNFTerm.blockCode_positions, FABL.DNFTerm.coordinatesAtPositions, FABL.DNFTerm.width, FABL.SwitchingBlockCode.termCoordinates}
\difficulty{4}
Let $T$ be a DNF term of width at most $w$, let $S$ be a set of positions among its
literals, and let $x$ be a binary assignment on $n$ coordinates. The list of term
coordinates selected by the block code of $T$ formed from $S$ and $x$ is exactly the
list of coordinates of the literals in $S$.
\end{theorem}

\begin{definition}[Assignment for a Term and Block Code]
\lean{FABL.SwitchingBlockCode.assignmentForTerm}
\label{FABL.SwitchingBlockCode.assignmentForTerm}
\leanok
\uses{FABL.DNFTerm, FABL.DNFTerm.positionOfSupport, FABL.DNFTerm.requiredAt, FABL.DNFTerm.support, FABL.SwitchingBlockCode, FABL.binarySignEquiv}
Let $B$ be a switching block code of length $w$ and let $T$ be a DNF term on $n$
variables. On a supported coordinate of $T$, use the sum modulo $2$ of the bit of $B$ at
that coordinate’s literal position and the bit encoding the literal’s required sign,
provided that the position is less than $w$; use zero otherwise. Coordinates outside the
support of $T$ are also assigned zero.
\end{definition}

\begin{theorem}[Agreement of term assignments]
\lean{FABL.SwitchingBlockCode.assignmentForTerm_eq_assignmentOfBlock}
\label{FABL.SwitchingBlockCode.assignmentForTerm_eq_assignmentOfBlock}
\leanok
\uses{FABL.DNFTerm, FABL.DNFTerm.assignmentOfBlock, FABL.DNFTerm.positionOfSupport, FABL.DNFTerm.support, FABL.DNFTerm.width, FABL.SwitchingBlockCode, FABL.SwitchingBlockCode.assignmentForTerm}
\difficulty{4}
Let $T$ be a conjunction of literals on $n$ variables, with distinct variable indices,
and let its width be at most $w$. For any width-$w$ switching block code $B$, the
assignment obtained by placing the code bits at the term’s literal positions and
adjusting each by the bit representing that literal’s required sign, with zero assigned
outside the term’s support, is the same assignment as the one defined directly from $T$
and $B$.
\end{theorem}

\begin{theorem}[Assignment from a block code]
\lean{FABL.DNFTerm.assignmentForTerm_blockCode_literal}
\label{FABL.DNFTerm.assignmentForTerm_blockCode_literal}
\leanok
\uses{FABL.DNFTerm, FABL.DNFTerm.assignmentOfBlock_blockCode_literal, FABL.DNFTerm.blockCode, FABL.DNFTerm.width, FABL.SwitchingBlockCode.assignmentForTerm, FABL.SwitchingBlockCode.assignmentForTerm_eq_assignmentOfBlock}
\difficulty{3}
Let $T$ be a conjunction of literals on distinct coordinates, with width at most $w$,
and let $P$ be a set of positions among its literals. Given an assignment
$x\in\mathbb{F}_2^n$, form the block code on the embedded positions using, at each such
position, the corresponding assignment bit adjusted by the literal's required sign. The
assignment constructed from this block code for $T$ agrees with $x$ at the coordinate of
every literal whose position lies in $P$.
\end{theorem}

\begin{definition}[Satisfying Assignment for a DNF Term]
\lean{FABL.DNFTerm.satisfyingAssignment}
\label{FABL.DNFTerm.satisfyingAssignment}
\leanok
\uses{FABL.DNFTerm, FABL.DNFTerm.requiredAt, FABL.DNFTerm.support, FABL.binarySignEquiv}
Let $T$ be a DNF term on $n$ variables. The assignment gives each supported coordinate
the bit encoding the sign required by its literal, and gives every coordinate outside
the support the value zero.
\end{definition}

\begin{theorem}[The satisfying assignment meets every literal]
\lean{FABL.DNFTerm.satisfyingAssignment_literal}
\label{FABL.DNFTerm.satisfyingAssignment_literal}
\leanok
\uses{FABL.DNFTerm, FABL.DNFTerm.literalAt_eq, FABL.DNFTerm.mem_support_of_mem_literals, FABL.DNFTerm.requiredAt, FABL.DNFTerm.satisfyingAssignment, FABL.binarySignEquiv}
\difficulty{4}
Let $T$ be a DNF term on $n$ coordinates, and let $j$ be a position in its list of
literals. The satisfying assignment of $T$, evaluated at the coordinate indexed by the
literal in position $j$, equals the binary value corresponding to that literal's
required sign.
\end{theorem}

\begin{theorem}[Relative bits of a satisfying assignment]
\lean{FABL.DNFTerm.relativeBit_satisfyingAssignment}
\label{FABL.DNFTerm.relativeBit_satisfyingAssignment}
\leanok
\uses{FABL.DNFTerm, FABL.DNFTerm.relativeBit, FABL.DNFTerm.satisfyingAssignment, FABL.DNFTerm.satisfyingAssignment_literal}
\difficulty{3}
Let $n$ be a natural number and let $T$ be a DNF term on $n$ coordinates. Define an
assignment in $oldsymbol{F}_2^n$ by setting each coordinate in the support of $T$ to
the bit corresponding to the required sign of its literal, and each coordinate outside
the support to $0$. Then, at every literal position in $T$, the assignment's coordinate
plus the bit corresponding to that literal's required sign is $0$ in $oldsymbol{F}_2$.
\end{theorem}

\begin{definition}[Satisfying extension]
\lean{FABL.DNFTerm.satisfyingExtension}
\label{FABL.DNFTerm.satisfyingExtension}
\leanok
\uses{FABL.CoordRestriction, FABL.CoordRestriction.fixCoordinateList, FABL.DNFTerm, FABL.DNFTerm.coordinatesAtPositions, FABL.DNFTerm.satisfyingAssignment}
Given a DNF term $T$ on $n$ coordinates, a coordinate restriction $\rho$, and a set of
positions among the literals of $T$, the satisfying extension is the restriction
obtained by fixing the coordinates selected at those positions according to a satisfying
assignment of $T$.
\end{definition}

\begin{theorem}[Compatibility after fixing selected literals]
\lean{FABL.DNFTerm.compatible_fix_positions}
\label{FABL.DNFTerm.compatible_fix_positions}
\leanok
\uses{FABL.CoordRestriction, FABL.CoordRestriction.complete_fixCoordinateList, FABL.CoordRestriction.fixCoordinateList, FABL.DNFTerm, FABL.DNFTerm.Compatible, FABL.DNFTerm.coordinatesAtPositions, FABL.DNFTerm.coordinatesAtPositions_nodup, FABL.DNFTerm.eval_eq_neg_one_iff, FABL.DNFTerm.freeLiteralPositions, FABL.DNFTerm.mem_coordinatesAtPositions_iff, FABL.DNFTerm.mem_freeLiteralPositions_iff, FABL.DNFTerm.relativeBit, FABL.binarySignEquiv, FABL.signEncode}
\difficulty{5}
Let $T$ be a conjunction of literals on $n$ Boolean coordinates, let $\rho$ be a
coordinate restriction compatible with $T$, and let $S$ be a set of literal positions
whose coordinates are free under $\rho$. For $x\in\mathbb{F}_2^n$, suppose that the
relative bit of $x$ is $0$ at every free literal position. Fix the coordinates of the
literals in $S$ to the signs encoded by $x$. The resulting coordinate restriction
remains compatible with $T$.
\end{theorem}

\begin{theorem}[Compatibility after fixing free literals]
\lean{FABL.DNFTerm.compatible_fix_freeLiteralPositions}
\label{FABL.DNFTerm.compatible_fix_freeLiteralPositions}
\leanok
\uses{FABL.CoordRestriction, FABL.CoordRestriction.complete_fixCoordinateList, FABL.CoordRestriction.fixCoordinateList, FABL.DNFTerm, FABL.DNFTerm.Compatible, FABL.DNFTerm.coordinatesAtPositions, FABL.DNFTerm.coordinatesAtPositions_nodup, FABL.DNFTerm.coordinatesAt_freeLiteralPositions_all_free, FABL.DNFTerm.eval_eq_neg_one_iff, FABL.DNFTerm.freeLiteralPositions, FABL.DNFTerm.mem_freeLiteralPositions_iff, FABL.DNFTerm.relativeBit, FABL.binarySignEquiv, FABL.signEncode}
\difficulty{5}
Let $T$ be a conjunction of literals on $n$ Boolean coordinates, and let $\rho$ be a
coordinate restriction compatible with $T$. For $x\in\mathbb{F}_2^n$, suppose that at
every position corresponding to a literal whose coordinate is free under $\rho$, the
relative bit of $x$ is $0$. Fix all such free literal coordinates to the signs encoded
by $x$. The resulting coordinate restriction remains compatible with $T$.
\end{theorem}

\begin{theorem}[Nonzero block code from incompatibility]
\lean{FABL.DNFTerm.blockCode_nonzero_of_not_compatible}
\label{FABL.DNFTerm.blockCode_nonzero_of_not_compatible}
\leanok
\uses{FABL.CoordRestriction, FABL.CoordRestriction.fixCoordinateList, FABL.DNFTerm, FABL.DNFTerm.Compatible, FABL.DNFTerm.blockCode, FABL.DNFTerm.blockCode_bits, FABL.DNFTerm.blockCode_positions, FABL.DNFTerm.compatible_fix_freeLiteralPositions, FABL.DNFTerm.coordinatesAtPositions, FABL.DNFTerm.freeLiteralPositions, FABL.DNFTerm.positionEmbedding, FABL.DNFTerm.relativeBit, FABL.DNFTerm.width, FABL.SwitchingBlockCode.Nonzero}
\difficulty{4}
Let $T$ be a conjunction of literals on $n$ Boolean coordinates, let its width be at
most $w$, and let $\rho$ be a coordinate restriction. Suppose some Boolean assignment
makes every literal of $T$ true while respecting $\rho$. For $x\in\mathbb{F}_2^n$, fix
each coordinate that is free under $\rho$ and appears in $T$ to the sign encoded by the
corresponding coordinate of $x$. If the resulting restriction is incompatible with $T$,
then the block code of these free literal positions, formed using $x$, has at least one
position whose bit is $1$.
\end{theorem}

\begin{theorem}[Satisfying extension preserves compatibility]
\lean{FABL.DNFTerm.satisfyingExtension_compatible}
\label{FABL.DNFTerm.satisfyingExtension_compatible}
\leanok
\uses{FABL.CoordRestriction, FABL.DNFTerm, FABL.DNFTerm.Compatible, FABL.DNFTerm.compatible_fix_freeLiteralPositions, FABL.DNFTerm.freeLiteralPositions, FABL.DNFTerm.relativeBit_satisfyingAssignment, FABL.DNFTerm.satisfyingAssignment, FABL.DNFTerm.satisfyingExtension}
\difficulty{3}
Let $T$ be a DNF term on $n$ coordinates and let $
ho$ be a coordinate restriction
compatible with $T$, meaning that some completion of its free coordinates makes $T$
evaluate to $-1$. Then the restriction obtained by fixing the coordinates at the free
literal positions to the values prescribed by the satisfying assignment of $T$ is also
compatible with $T$.
\end{theorem}

\begin{theorem}[Depth in a queried-coordinate tree]
\lean{FABL.CoordinateDecisionTree.exists_finish_depth_ge_sub_of_depth_queryCoordinates}
\label{FABL.CoordinateDecisionTree.exists_finish_depth_ge_sub_of_depth_queryCoordinates}
\leanok
\uses{FABL.CoordinateDecisionTree, FABL.CoordinateDecisionTree.depth, FABL.CoordinateDecisionTree.queryCoordinates}
\difficulty{5}
Let $C$ be a list of coordinates, $s\in\mathbb{F}_2^n$ a seed, and
$F:\mathbb{F}_2^n\to\mathcal{T}$ a map into coordinate decision trees. If the tree
formed by querying the coordinates in $C$, branching over both bit values and updating
the seed accordingly, has depth at least $k$, and $|C|<k$, then some
$x\in\mathbb{F}_2^n$ satisfies $\operatorname{depth}(F(x))\ge k-|C|$.
\end{theorem}

\begin{definition}[Canonical compatible term]
\lean{FABL.DNFFormula.canonicalTerm}
\label{FABL.DNFFormula.canonicalTerm}
\leanok
\uses{FABL.CoordRestriction, FABL.DNFFormula, FABL.DNFFormula.canonicalTermIndex, FABL.DNFFormula.compatibleTermIndices, FABL.DNFTerm}
Given a DNF formula $\varphi$ on $n$ coordinates, a coordinate restriction $\rho$, and
the assumption that at least one term of $\varphi$ is compatible with $\rho$, the
canonical term is the term at the index selected by the formula's canonical-term rule
for $\rho$.
\end{definition}

\begin{theorem}[Membership of the canonical term]
\lean{FABL.DNFFormula.canonicalTerm_mem}
\label{FABL.DNFFormula.canonicalTerm_mem}
\leanok
\uses{FABL.CoordRestriction, FABL.DNFFormula, FABL.DNFFormula.canonicalTerm, FABL.DNFFormula.canonicalTermIndex, FABL.DNFFormula.compatibleTermIndices}
\difficulty{1}
Let $\phi$ be a DNF formula on $n$ coordinates and let $\rho$ be a coordinate
restriction. If at least one term is compatible with $\rho$, then the canonical term
selected for $\rho$ belongs to the list of terms of $\phi$.
\end{theorem}

\begin{theorem}[Width bound for the canonical term]
\lean{FABL.DNFFormula.canonicalTerm_width_le}
\label{FABL.DNFFormula.canonicalTerm_width_le}
\leanok
\uses{FABL.CoordRestriction, FABL.DNFFormula, FABL.DNFFormula.canonicalTerm, FABL.DNFFormula.canonicalTerm_mem, FABL.DNFFormula.compatibleTermIndices, FABL.DNFFormula.width, FABL.DNFFormula.width_le_of_mem, FABL.DNFTerm.width}
\difficulty{1}
Let $\phi$ be a DNF formula on $n$ coordinates, and suppose its width is at most $w$.
For any coordinate restriction $\rho$ under which at least one term is compatible, the
canonical term selected for $\rho$ has width at most $w$.
\end{theorem}

\begin{theorem}[Compatibility of the canonical term]
\lean{FABL.DNFFormula.canonicalTerm_compatible'}
\label{FABL.DNFFormula.canonicalTerm_compatible'}
\leanok
\uses{FABL.CoordRestriction, FABL.DNFFormula, FABL.DNFFormula.canonicalTerm, FABL.DNFFormula.canonicalTerm_compatible, FABL.DNFFormula.compatibleTermIndices, FABL.DNFTerm.Compatible}
\difficulty{1}
Let $\phi$ be a DNF formula on $n$ coordinates and let $\rho$ be a coordinate
restriction. If at least one term of $\phi$ is compatible with $\rho$, then the
canonical term selected by least index is compatible with $\rho$.
\end{theorem}

\begin{definition}[Switching traces for DNF formulas]
\lean{FABL.DNFFormula.SwitchingTrace}
\statementsource{arora-barak-ch13-circuit-lowerbounds}{
  pdf-52929dbaa962-p003-b004,
  pdf-52929dbaa962-p003-b005,
  pdf-52929dbaa962-p004-b001,
  pdf-52929dbaa962-p004-b002,
  pdf-52929dbaa962-p005-b001
}
\label{FABL.DNFFormula.SwitchingTrace}
\leanok
\uses{FABL.CoordRestriction, FABL.CoordRestriction.fixCoordinateList, FABL.DNFFormula, FABL.DNFFormula.canonicalTerm, FABL.DNFFormula.compatibleTermIndices, FABL.DNFTerm.Compatible, FABL.DNFTerm.coordinatesAtPositions, FABL.DNFTerm.freeLiteralPositions}
Let $\varphi$ be a DNF formula on $n$ coordinates, let $\rho$ assign each coordinate one
of the states free, fixed at $1$, or fixed at $-1$, and let $k\in\mathbb{N}$. A
switching trace of index $k+1$ is either specified by a compatible term chosen
canonically from $\varphi$ under $\rho$ and a set of $k+1$ of its free literal
positions, or specified by a compatible canonical term under $\rho$, an index $i<k$, and
a vector $x\in\mathbb{F}_2^n$ such that the term has exactly $i+1$ free literal
positions and becomes incompatible after the coordinates at all those positions are
fixed according to $x$; in the latter case the trace also includes a switching trace of
index $k-i$ under the resulting coordinate restriction.
\end{definition}

\begin{theorem}[Switching trace from canonical block tree depth]
\lean{FABL.DNFFormula.nonempty_switchingTrace_of_depth_canonicalBlockTreeAux}
\label{FABL.DNFFormula.nonempty_switchingTrace_of_depth_canonicalBlockTreeAux}
\leanok
\uses{FABL.CoordRestriction, FABL.CoordRestriction.fixCoordinateList, FABL.CoordinateDecisionTree.depth, FABL.CoordinateDecisionTree.exists_finish_depth_ge_sub_of_depth_queryCoordinates, FABL.DNFFormula, FABL.DNFFormula.SwitchingTrace, FABL.DNFFormula.canonicalBlockTreeAux, FABL.DNFFormula.canonicalTerm, FABL.DNFFormula.canonicalTermIndex, FABL.DNFFormula.compatibleTermIndices, FABL.DNFTerm.Compatible, FABL.DNFTerm.coordinatesAtPositions, FABL.DNFTerm.freeLiteralPositions}
\difficulty{6}
Let $n$ be a natural number, let $\varphi$ be a DNF formula on $n$ coordinates, let $m
\in \mathbb{N}$ be a fuel parameter, let $\rho$ be a coordinate restriction, and let $k
\in \mathbb{N}$. If the depth of the canonical block decision tree constructed for
$\varphi$, with fuel $m$ and restriction $\rho$, is at least $k+1$, then there exists a
switching trace for $\varphi$ relative to $\rho$ of parameter $k+1$.
\end{theorem}

\begin{theorem}[Switching trace from decision tree depth]
\lean{FABL.DNFFormula.nonempty_switchingTrace_of_decisionTreeDepth}
\label{FABL.DNFFormula.nonempty_switchingTrace_of_decisionTreeDepth}
\leanok
\uses{FABL.CoordRestriction, FABL.CoordRestriction.restrict, FABL.CoordinateDecisionTree.depth, FABL.DNFFormula, FABL.DNFFormula.SwitchingTrace, FABL.DNFFormula.canonicalBlockTree, FABL.DNFFormula.canonicalBlockTreeAux, FABL.DNFFormula.canonicalBlockTreeRaw, FABL.DNFFormula.decisionTreeDepth_le_canonicalBlockTree_depth, FABL.DNFFormula.nonempty_switchingTrace_of_depth_canonicalBlockTreeAux, FABL.DNFFormula.toBooleanFunction, FABL.F₂DecisionTree.decisionTreeDepth, FABL.F₂DecisionTree.depth}
\difficulty{4}
Let $n$ be a natural number, let $\varphi$ be a DNF formula on $n$ coordinates, let
$\rho$ be a coordinate restriction, and let $k \in \mathbb{N}$. If the decision-tree
depth of the restriction of the Boolean function represented by $\varphi$ under $\rho$
is at least $k+1$, then there exists a switching trace for $\varphi$ relative to $\rho$
of parameter $k+1$.
\end{theorem}

\begin{definition}[Target Restriction of a Switching Trace]
\lean{FABL.DNFFormula.SwitchingTrace.target}
\statementsource{arora-barak-ch13-circuit-lowerbounds}{
  pdf-52929dbaa962-p003-b004,
  pdf-52929dbaa962-p003-b005,
  pdf-52929dbaa962-p004-b001,
  pdf-52929dbaa962-p004-b002,
  pdf-52929dbaa962-p005-b001
}
\label{FABL.DNFFormula.SwitchingTrace.target}
\leanok
\uses{FABL.CoordRestriction, FABL.DNFFormula.SwitchingTrace, FABL.DNFFormula.canonicalTerm, FABL.DNFTerm.freeLiteralPositions, FABL.DNFTerm.satisfyingExtension}
Let $\varphi$ be a DNF formula on $n$ coordinates and let a switching trace for
$\varphi$ start at restriction $\rho$. Its target is a coordinate-state assignment: for
a final trace, it is the satisfying extension of the canonical compatible term using the
chosen positions; for a recursive trace, it is the satisfying extension of the canonical
compatible term at the current restriction, applied to the tail's target and using all
currently free literal positions.
\end{definition}

\begin{definition}[Switching Code of a Trace]
\lean{FABL.DNFFormula.SwitchingTrace.code}
\statementsource{arora-barak-ch13-circuit-lowerbounds}{
  pdf-52929dbaa962-p003-b004,
  pdf-52929dbaa962-p003-b005,
  pdf-52929dbaa962-p004-b001,
  pdf-52929dbaa962-p004-b002,
  pdf-52929dbaa962-p005-b001
}
\label{FABL.DNFFormula.SwitchingTrace.code}
\leanok
\uses{FABL.CoordRestriction, FABL.DNFFormula.SwitchingTrace, FABL.DNFFormula.canonicalTerm, FABL.DNFFormula.canonicalTerm_compatible', FABL.DNFFormula.canonicalTerm_width_le, FABL.DNFFormula.width, FABL.DNFTerm.blockCode, FABL.DNFTerm.blockCode_canonical, FABL.DNFTerm.blockCode_card_positions, FABL.DNFTerm.blockCode_nonzero_of_not_compatible, FABL.DNFTerm.freeLiteralPositions, FABL.DNFTerm.satisfyingAssignment, FABL.SwitchingCode}
Let $\varphi$ be a DNF formula of width at most $w$, and let a switching trace for
$\varphi$ have index $k$. Its code is a switching code of width $w$ and index $k$. A
final trace yields a final block code from the chosen literal positions and a satisfying
assignment of the canonical compatible term; a recursive trace yields a nonzero block
code from the current free literal positions and the chosen binary vector, followed by
the code of the tail trace.
\end{definition}

\begin{definition}[Decoding a Switching Code]
\lean{FABL.DNFFormula.decodeSwitchingCode}
\statementsource{arora-barak-ch13-circuit-lowerbounds}{
  pdf-52929dbaa962-p004-b001,
  pdf-52929dbaa962-p004-b002,
  pdf-52929dbaa962-p005-b001
}
\label{FABL.DNFFormula.decodeSwitchingCode}
\leanok
\uses{FABL.CoordRestriction, FABL.CoordRestriction.fixCoordinateList, FABL.CoordRestriction.freeCoordinateList, FABL.DNFFormula, FABL.DNFFormula.canonicalTerm, FABL.DNFFormula.compatibleTermIndices, FABL.DNFFormula.width, FABL.SwitchingBlockCode.assignmentForTerm, FABL.SwitchingBlockCode.termCoordinates, FABL.SwitchingCode}
Let $\varphi$ be a DNF formula of width at most $w$, let $\rho$ be a coordinate-state
assignment on $n$ coordinates, and let $C$ be a switching code of width $w$ and index
$k$. The decoded assignment is defined recursively: if no term is compatible with the
current restriction, it leaves that restriction unchanged. Otherwise, for a final code
it makes free the coordinates selected by the code in the canonical compatible term; for
a recursive code it first fixes those coordinates according to the assignment determined
by the block code and term, decodes the tail from that updated restriction, and then
makes those coordinates free.
\end{definition}

\begin{theorem}[Canonical term under an extending restriction]
\lean{FABL.DNFFormula.canonicalTerm_eq_of_extends}
\label{FABL.DNFFormula.canonicalTerm_eq_of_extends}
\leanok
\uses{FABL.CoordRestriction, FABL.CoordRestriction.Extends, FABL.DNFFormula, FABL.DNFFormula.canonicalTerm, FABL.DNFFormula.canonicalTermIndex_eq_of_extends, FABL.DNFFormula.compatibleTermIndices, FABL.DNFTerm.Compatible}
\difficulty{3}
Let $\phi$ be a DNF formula on $n$ coordinates, and let $\rho,\eta$ be coordinate
restrictions such that $\eta$ agrees with $\rho$ wherever $\rho$ is fixed. Suppose each
restriction admits at least one compatible term, and the canonical term selected for
$\rho$ is compatible with $\eta$. Then the canonical terms selected for $\rho$ and
$\eta$ are equal.
\end{theorem}

\begin{theorem}[Target extension and term compatibility]
\lean{FABL.DNFFormula.SwitchingTrace.target_extends_and_canonical_compatible}
\label{FABL.DNFFormula.SwitchingTrace.target_extends_and_canonical_compatible}
\leanok
\uses{FABL.CoordRestriction.Extends, FABL.CoordRestriction.Extends.parallel_fixCoordinateList, FABL.CoordRestriction.Extends.trans, FABL.CoordRestriction.extends_fixCoordinateList, FABL.CoordRestriction.fixCoordinateList, FABL.DNFFormula.SwitchingTrace, FABL.DNFFormula.SwitchingTrace.target, FABL.DNFFormula.canonicalTerm, FABL.DNFFormula.canonicalTerm_compatible', FABL.DNFFormula.compatibleTermIndices, FABL.DNFTerm.Compatible, FABL.DNFTerm.compatible_fix_positions, FABL.DNFTerm.compatible_of_extends_of_no_free, FABL.DNFTerm.coordinatesAtPositions, FABL.DNFTerm.coordinatesAtPositions_nodup, FABL.DNFTerm.coordinatesAt_freeLiteralPositions_all_free, FABL.DNFTerm.freeLiteralPositions, FABL.DNFTerm.freeLiteralPositions_fixCoordinateList_eq_empty, FABL.DNFTerm.mem_coordinatesAtPositions_iff, FABL.DNFTerm.mem_freeLiteralPositions_iff, FABL.DNFTerm.relativeBit_satisfyingAssignment, FABL.DNFTerm.satisfyingAssignment, FABL.DNFTerm.satisfyingExtension, FABL.DNFTerm.satisfyingExtension_compatible}
\difficulty{5}
Let $n$ be a natural number, let $\varphi$ be a disjunctive normal form on $n$ sign
variables, and let $\rho$ be a coordinate restriction. For every switching trace of
$\varphi$ starting from $\rho$, its target restriction extends $\rho$ by preserving
every coordinate that $\rho$ fixes, and there is a term of $\varphi$ compatible with
$\rho$ that is also compatible with the target restriction.
\end{theorem}

\begin{theorem}[Switching trace target extension]
\lean{FABL.DNFFormula.SwitchingTrace.target_extends}
\proofsource{arora-barak-ch13-circuit-lowerbounds}{
  pdf-52929dbaa962-p004-b002,
  pdf-52929dbaa962-p005-b001
}
\label{FABL.DNFFormula.SwitchingTrace.target_extends}
\leanok
\uses{FABL.CoordRestriction.Extends, FABL.DNFFormula.SwitchingTrace, FABL.DNFFormula.SwitchingTrace.target, FABL.DNFFormula.SwitchingTrace.target_extends_and_canonical_compatible}
\difficulty{1}
Let $n$ be a natural number, let $\varphi$ be a disjunctive normal form on $n$ sign
variables, and let $\rho$ be a coordinate restriction. For every switching trace of
$\varphi$ starting from $\rho$, its target restriction agrees with $\rho$ at every
coordinate that $\rho$ fixes.
\end{theorem}

\begin{theorem}[Compatible canonical term at the target]
\lean{FABL.DNFFormula.SwitchingTrace.exists_canonical_compatible_target}
\label{FABL.DNFFormula.SwitchingTrace.exists_canonical_compatible_target}
\leanok
\uses{FABL.DNFFormula.SwitchingTrace, FABL.DNFFormula.SwitchingTrace.target, FABL.DNFFormula.SwitchingTrace.target_extends_and_canonical_compatible, FABL.DNFFormula.canonicalTerm, FABL.DNFFormula.compatibleTermIndices, FABL.DNFTerm.Compatible}
\difficulty{1}
Let $\phi$ be a DNF formula, let $\rho$ be a coordinate restriction, and let
$\mathrm{tr}$ be a switching trace for $\phi$ based at $\rho$. There is a compatible
term of $\phi$ that is canonical for some nonempty set of compatible term indices at
$\rho$ and is itself compatible with the target restriction associated with
$\mathrm{tr}$.
\end{theorem}

\begin{theorem}[Decoding a switching trace]
\lean{FABL.DNFFormula.SwitchingTrace.decode_target_code}
\proofsource{arora-barak-ch13-circuit-lowerbounds}{
  pdf-52929dbaa962-p004-b001,
  pdf-52929dbaa962-p004-b002,
  pdf-52929dbaa962-p005-b001
}
\label{FABL.DNFFormula.SwitchingTrace.decode_target_code}
\leanok
\uses{FABL.CoordRestriction.Extends, FABL.CoordRestriction.Extends.parallel_fixCoordinateList, FABL.CoordRestriction.Extends.trans, FABL.CoordRestriction.extends_fixCoordinateList, FABL.CoordRestriction.fixCoordinateList, FABL.CoordRestriction.fixCoordinateList_apply_of_mem, FABL.CoordRestriction.fixCoordinateList_eq_self, FABL.CoordRestriction.fixCoordinateList_overwrite, FABL.CoordRestriction.fixedState, FABL.CoordRestriction.fixedState_ne_free, FABL.CoordRestriction.freeCoordinateList, FABL.CoordRestriction.freeCoordinateList_fixCoordinateList, FABL.DNFFormula.SwitchingTrace, FABL.DNFFormula.SwitchingTrace.code, FABL.DNFFormula.SwitchingTrace.target, FABL.DNFFormula.SwitchingTrace.target_extends, FABL.DNFFormula.canonicalTerm, FABL.DNFFormula.canonicalTermIndex, FABL.DNFFormula.canonicalTerm_compatible', FABL.DNFFormula.canonicalTerm_eq_of_extends, FABL.DNFFormula.canonicalTerm_width_le, FABL.DNFFormula.compatibleTermIndices, FABL.DNFFormula.decodeSwitchingCode, FABL.DNFFormula.width, FABL.DNFTerm.Compatible, FABL.DNFTerm.assignmentOfBlock, FABL.DNFTerm.assignmentOfBlock_blockCode_literal, FABL.DNFTerm.blockCode, FABL.DNFTerm.compatible_fix_positions, FABL.DNFTerm.compatible_of_extends_of_no_free, FABL.DNFTerm.coordinatesAtPositions, FABL.DNFTerm.coordinatesAtPositions_nodup, FABL.DNFTerm.coordinatesAt_freeLiteralPositions_all_free, FABL.DNFTerm.freeLiteralPositions, FABL.DNFTerm.freeLiteralPositions_fixCoordinateList_eq_empty, FABL.DNFTerm.mem_coordinatesAtPositions_iff, FABL.DNFTerm.mem_freeLiteralPositions_iff, FABL.DNFTerm.relativeBit_satisfyingAssignment, FABL.DNFTerm.satisfyingAssignment, FABL.DNFTerm.satisfyingExtension, FABL.DNFTerm.satisfyingExtension_compatible, FABL.DNFTerm.termCoordinates_blockCode, FABL.DNFTerm.width, FABL.SwitchingBlockCode.assignmentForTerm, FABL.SwitchingBlockCode.assignmentForTerm_eq_assignmentOfBlock, FABL.SwitchingBlockCode.termCoordinates, FABL.signEncode}
\difficulty{6}
Let $n,w\in\mathbb{N}$, let $\phi$ be a DNF formula on $n$ coordinates with width at
most $w$, and let $\rho$ be a restriction assigning each coordinate one of the states
free, fixed to $1$, or fixed to $-1$. For every switching trace $\mathrm{tr}$ of $\phi$
based at $\rho$, decoding the code associated with $\mathrm{tr}$ at the target
restriction of $\mathrm{tr}$ recovers $\rho$.
\end{theorem}

\begin{definition}[Switching atom ratio]
\lean{FABL.switchingAtomRatio}
\label{FABL.switchingAtomRatio}
\leanok
For a real number $\delta$, the switching atom ratio is $\frac{2\delta}{1-\delta}$.
\end{definition}

\begin{theorem}[Weight change when fixing a coordinate]
\lean{FABL.restrictionAssignmentWeightAt_fixCoordinate}
\label{FABL.restrictionAssignmentWeightAt_fixCoordinate}
\leanok
\uses{FABL.CoordRestriction, FABL.CoordRestriction.fixCoordinate, FABL.CoordRestriction.fixCoordinate_apply_of_ne, FABL.CoordRestriction.fixCoordinate_apply_self, FABL.Sign, FABL.coordRestrictionWeightAt, FABL.coordRestrictionWeightAt_fixedState, FABL.restrictionAssignmentWeightAt, FABL.switchingAtomRatio}
\difficulty{4}
Let $n$ be a natural number, let $\delta<1$, and let $\rho$ be an assignment of a
restriction state to each coordinate in $[n]$. If coordinate $i$ is free in $\rho$, then
for either sign $s\in\{\pm1\}$, the weight of $\rho$ equals $2\delta/(1-\delta)$ times
the weight obtained by fixing coordinate $i$ to $s$ and leaving every other restriction
state unchanged.
\end{theorem}

\begin{theorem}[Weight change for a list of fixed coordinates]
\lean{FABL.restrictionAssignmentWeightAt_fixCoordinateList}
\label{FABL.restrictionAssignmentWeightAt_fixCoordinateList}
\leanok
\uses{FABL.CoordRestriction, FABL.CoordRestriction.fixCoordinate, FABL.CoordRestriction.fixCoordinateList, FABL.CoordRestriction.fixCoordinateList_cons, FABL.CoordRestriction.fixCoordinate_apply_of_ne, FABL.restrictionAssignmentWeightAt, FABL.restrictionAssignmentWeightAt_fixCoordinate, FABL.signEncode, FABL.switchingAtomRatio}
\difficulty{5}
Let $n$ be a natural number, let $\delta<1$, let $\rho$ assign a restriction state to
each coordinate in $[n]$, let $L$ be a list of distinct coordinates, and let $x$ be a
binary vector of length $n$. If every coordinate in $L$ is free in $\rho$, then the
weight of $\rho$ equals $\bigl(2\delta/(1-\delta)\bigr)^{|L|}$ times the weight obtained
by fixing each coordinate in $L$ according to the corresponding sign encoded by $x$.
\end{theorem}

\begin{theorem}[Weight invariance on already fixed coordinates]
\lean{FABL.restrictionAssignmentWeightAt_fixCoordinateList_of_all_fixed}
\label{FABL.restrictionAssignmentWeightAt_fixCoordinateList_of_all_fixed}
\leanok
\uses{FABL.CoordRestriction, FABL.CoordRestriction.fixCoordinateList, FABL.CoordRestriction.fixCoordinateList_apply_of_mem, FABL.CoordRestriction.fixCoordinateList_apply_of_not_mem, FABL.coordRestrictionWeightAt_fixedState, FABL.restrictionAssignmentWeightAt}
\difficulty{4}
Let $n$ be a natural number, let $\delta\in\mathbb R$, let $\rho$ assign a restriction
state to each coordinate in $[n]$, let $L$ be a list of distinct coordinates, and let
$x$ be a binary vector of length $n$. If every coordinate in $L$ is already fixed in
$\rho$, then fixing those coordinates according to the signs encoded by $x$ leaves the
weight of $\rho$ unchanged.
\end{theorem}

\begin{theorem}[Restriction weight along a switching trace]
\lean{FABL.DNFFormula.SwitchingTrace.restrictionAssignmentWeightAt_eq_ratio_pow_target}
\label{FABL.DNFFormula.SwitchingTrace.restrictionAssignmentWeightAt_eq_ratio_pow_target}
\leanok
\uses{FABL.CoordRestriction.Extends, FABL.CoordRestriction.fixCoordinateList, FABL.CoordRestriction.fixCoordinateList_apply_of_mem, FABL.CoordRestriction.fixedState_ne_free, FABL.DNFFormula.SwitchingTrace, FABL.DNFFormula.SwitchingTrace.target, FABL.DNFFormula.SwitchingTrace.target_extends, FABL.DNFFormula.canonicalTerm, FABL.DNFTerm.coordinatesAtPositions, FABL.DNFTerm.coordinatesAtPositions_nodup, FABL.DNFTerm.coordinatesAt_freeLiteralPositions_all_free, FABL.DNFTerm.freeLiteralPositions, FABL.DNFTerm.mem_coordinatesAtPositions_iff, FABL.DNFTerm.mem_freeLiteralPositions_iff, FABL.DNFTerm.satisfyingAssignment, FABL.DNFTerm.satisfyingExtension, FABL.restrictionAssignmentWeightAt, FABL.restrictionAssignmentWeightAt_fixCoordinateList, FABL.restrictionAssignmentWeightAt_fixCoordinateList_of_all_fixed, FABL.switchingAtomRatio}
\difficulty{6}
Let $\phi$ be a DNF formula, let $\rho$ be a coordinate restriction, and let
$\mathrm{tr}$ be a switching trace of length $k$ based at $\rho$. For $\delta<1$, assign
weight $\delta$ to each free coordinate and weight $(1-\delta)/2$ to each fixed
coordinate, and let the weight of a restriction be the product of its coordinate
weights. Then the weight of $\rho$ equals $\left(\frac{2\delta}{1-\delta}\right)^k$
times the weight of the trace's target restriction.
\end{theorem}

\begin{definition}[Bad Restrictions for a DNF Formula]
\lean{FABL.DNFFormula.SwitchingBad}
\statementsource{arora-barak-ch13-circuit-lowerbounds}{pdf-52929dbaa962-p003-b004}
\label{FABL.DNFFormula.SwitchingBad}
\leanok
\uses{FABL.CoordRestriction, FABL.CoordRestriction.restrict, FABL.DNFFormula, FABL.DNFFormula.toBooleanFunction, FABL.F₂DecisionTree.decisionTreeDepth}
Let $\varphi$ be a DNF formula on $n$ coordinates and let $k$ be a natural number. The
bad restrictions are the coordinate-state assignments $\rho$ for which the decision-tree
depth of the Boolean function represented by $\varphi$, after restriction by $\rho$, is
at least $k$.
\end{definition}

\begin{theorem}[Switching failure weight identity]
\lean{FABL.DNFFormula.sum_switchingBadWeight_eq_coordSwitchingFailureProbability}
\label{FABL.DNFFormula.sum_switchingBadWeight_eq_coordSwitchingFailureProbability}
\leanok
\uses{FABL.DNFFormula, FABL.DNFFormula.SwitchingBad, FABL.DNFFormula.toBooleanFunction, FABL.coordSwitchingFailureIndicator, FABL.coordSwitchingFailureProbability, FABL.restrictionAssignmentWeightAt}
\difficulty{3}
Let $n,k\in\mathbb{N}$, let $f$ be a Boolean function on the $n$-dimensional sign cube,
and let $\delta\in\mathbb{R}$. The sum, over all coordinate restrictions for which the
restricted function has decision-tree depth at least $k$, of the restriction weights
equals the coordinate switching failure probability of $f$ at parameters $\delta$ and
$k$.
\end{theorem}

\begin{definition}[Switching trace from a bad restriction]
\lean{FABL.DNFFormula.traceOfSwitchingBad}
\statementsource{arora-barak-ch13-circuit-lowerbounds}{
  pdf-52929dbaa962-p003-b004,
  pdf-52929dbaa962-p003-b005,
  pdf-52929dbaa962-p004-b001,
  pdf-52929dbaa962-p004-b002,
  pdf-52929dbaa962-p005-b001
}
\label{FABL.DNFFormula.traceOfSwitchingBad}
\leanok
\uses{FABL.DNFFormula, FABL.DNFFormula.SwitchingBad, FABL.DNFFormula.SwitchingTrace, FABL.DNFFormula.nonempty_switchingTrace_of_decisionTreeDepth}
Let $n,k$ be natural numbers, let $\varphi$ be a DNF formula on $n$ sign coordinates,
and let $\rho$ be a coordinate restriction whose restricted formula has decision-tree
depth at least $k+1$. The trace construction assigns to $\rho$ a switching trace of
index $k+1$.
\end{definition}

\begin{definition}[Switching encoding]
\lean{FABL.DNFFormula.switchingEncoding}
\statementsource{arora-barak-ch13-circuit-lowerbounds}{
  pdf-52929dbaa962-p003-b004,
  pdf-52929dbaa962-p003-b005,
  pdf-52929dbaa962-p004-b001,
  pdf-52929dbaa962-p004-b002,
  pdf-52929dbaa962-p005-b001
}
\label{FABL.DNFFormula.switchingEncoding}
\leanok
\uses{FABL.CoordRestriction, FABL.DNFFormula, FABL.DNFFormula.SwitchingBad, FABL.DNFFormula.SwitchingTrace.code, FABL.DNFFormula.SwitchingTrace.target, FABL.DNFFormula.traceOfSwitchingBad, FABL.DNFFormula.width, FABL.SwitchingCode}
Let $n,w,k$ be natural numbers and let $\varphi$ be a DNF formula on $n$ sign
coordinates whose width is at most $w$. The switching encoding maps each coordinate
restriction for which the restricted formula has decision-tree depth at least $k+1$ to a
pair consisting of a coordinate restriction and a switching code of width $w$ and index
$k+1$.
\end{definition}

\begin{theorem}[Injectivity of the switching encoding]
\lean{FABL.DNFFormula.switchingEncoding_injective}
\proofsource{arora-barak-ch13-circuit-lowerbounds}{
  pdf-52929dbaa962-p003-b004,
  pdf-52929dbaa962-p003-b005,
  pdf-52929dbaa962-p004-b001,
  pdf-52929dbaa962-p004-b002,
  pdf-52929dbaa962-p005-b001
}
\label{FABL.DNFFormula.switchingEncoding_injective}
\leanok
\uses{FABL.DNFFormula, FABL.DNFFormula.SwitchingTrace.code, FABL.DNFFormula.SwitchingTrace.decode_target_code, FABL.DNFFormula.SwitchingTrace.target, FABL.DNFFormula.decodeSwitchingCode, FABL.DNFFormula.switchingEncoding, FABL.DNFFormula.traceOfSwitchingBad, FABL.DNFFormula.width}
\difficulty{4}
Let $n,w,k\in\mathbb{N}$, and let $\varphi$ be a DNF formula on $n$ coordinates with
width at most $w$. The map sending each restriction for which the restricted formula has
decision-tree depth at least $k+1$ to its associated pair consisting of a target
restriction and a switching code of length $k+1$ is injective.
\end{theorem}

\begin{theorem}[Nonnegativity of the switching atom ratio]
\lean{FABL.switchingAtomRatio_nonneg}
\label{FABL.switchingAtomRatio_nonneg}
\leanok
\uses{FABL.switchingAtomRatio}
\difficulty{2}
For $\delta\in\mathbb R$ with $0\le\delta<1$, the quantity $\frac{2\delta}{1-\delta}$ is
nonnegative.
\end{theorem}

\begin{theorem}[Switching trace weight bound]
\lean{FABL.DNFFormula.sum_switchingTraceTargetWeight_le_card}
\proofsource{arora-barak-ch13-circuit-lowerbounds}{
  pdf-52929dbaa962-p003-b004,
  pdf-52929dbaa962-p003-b005,
  pdf-52929dbaa962-p005-b001
}
\label{FABL.DNFFormula.sum_switchingTraceTargetWeight_le_card}
\leanok
\uses{FABL.CoordRestriction, FABL.DNFFormula, FABL.DNFFormula.SwitchingBad, FABL.DNFFormula.SwitchingTrace.target, FABL.DNFFormula.switchingEncoding, FABL.DNFFormula.switchingEncoding_injective, FABL.DNFFormula.traceOfSwitchingBad, FABL.DNFFormula.width, FABL.SwitchingCode, FABL.restrictionAssignmentWeightAt, FABL.restrictionAssignmentWeightAt_nonneg, FABL.sum_restrictionAssignmentWeightAt}
\difficulty{4}
Let $n,w,k\in\mathbb{N}$, let $\varphi$ be a DNF formula on $n$ coordinates whose width
is at most $w$, and let $0\leq\delta\leq1$. Sum the weights, at parameter $\delta$, of
the target restrictions associated to the switching traces chosen for each restriction
on which the restricted formula has decision-tree depth at least $k+1$. This sum is at
most the number of switching codes of width parameter $w$ and length $k+1$.
\end{theorem}

\begin{theorem}[Coord switching failure probability le ratio pow card]
\lean{FABL.DNFFormula.coordSwitchingFailureProbability_le_ratio_pow_card}
\statementsource{arora-barak-ch13-circuit-lowerbounds}{
  pdf-52929dbaa962-p002-b003,
  pdf-52929dbaa962-p003-b004,
  pdf-52929dbaa962-p005-b001
}
\label{FABL.DNFFormula.coordSwitchingFailureProbability_le_ratio_pow_card}
\leanok
\uses{FABL.DNFFormula, FABL.DNFFormula.SwitchingBad, FABL.DNFFormula.SwitchingTrace.restrictionAssignmentWeightAt_eq_ratio_pow_target, FABL.DNFFormula.SwitchingTrace.target, FABL.DNFFormula.sum_switchingBadWeight_eq_coordSwitchingFailureProbability, FABL.DNFFormula.sum_switchingTraceTargetWeight_le_card, FABL.DNFFormula.toBooleanFunction, FABL.DNFFormula.traceOfSwitchingBad, FABL.DNFFormula.width, FABL.SwitchingCode, FABL.coordSwitchingFailureProbability, FABL.restrictionAssignmentWeightAt, FABL.switchingAtomRatio, FABL.switchingAtomRatio_nonneg}
\difficulty{4}
% dataset-pipeline: stub -- prose pending build_informal_statements.py
Statement pending regeneration.
\end{theorem}

\begin{theorem}[Failure probability at zero width]
\lean{FABL.DNFFormula.coordSwitchingFailureProbability_eq_zero_of_width_eq_zero}
\label{FABL.DNFFormula.coordSwitchingFailureProbability_eq_zero_of_width_eq_zero}
\leanok
\uses{FABL.CoordRestriction.complete, FABL.CoordRestriction.restrict, FABL.DNFFormula, FABL.DNFFormula.toBooleanFunction, FABL.DNFFormula.toBooleanFunction_eq_of_width_eq_zero, FABL.DNFFormula.width, FABL.F₂DecisionTree.decisionTreeDepth, FABL.coordSwitchingFailureIndicator, FABL.coordSwitchingFailureProbability, FABL.one_le_decisionTreeDepth_iff_nonconstant}
\difficulty{4}
Let $n$ be a nonnegative integer, let $\varphi$ be a disjunctive normal form on $n$ sign
coordinates whose width is zero, and let $\delta\in\mathbb R$ and $k\in\mathbb N$ with
$k>0$. The probability of coordinate-switching failure for $\varphi$ at parameter
$\delta$ and threshold $k$ is zero.
\end{theorem}

\begin{theorem}[Switching-code count at width one]
\lean{FABL.SwitchingCode.card_width_one}
\label{FABL.SwitchingCode.card_width_one}
\leanok
\uses{FABL.SwitchingCode, FABL.SwitchingCode.card_succ}
\difficulty{4}
For every positive integer $k$, there are exactly two switching codes of width one and
index $k$.
\end{theorem}

\begin{theorem}[Switching bound for DNF formulas]
\lean{FABL.hastadSwitchingLemma_dnf}
\proofsource{lmn-constant-depth-fourier}{pdf-63759bd909d5-p005-b004}
\proofsource{arora-barak-ch13-circuit-lowerbounds}{
  pdf-52929dbaa962-p002-b003,
  pdf-52929dbaa962-p003-b004,
  pdf-52929dbaa962-p004-b002,
  pdf-52929dbaa962-p005-b001
}
\proofsource{Analysis-of-Boolean-Functions-by-Ryan-ODonnell}{pdf-ffbe9c7b3306-p100-b005}
\label{FABL.hastadSwitchingLemma_dnf}
\leanok
\uses{FABL.DNFFormula, FABL.DNFFormula.coordSwitchingFailureProbability_eq_zero_of_width_eq_zero, FABL.DNFFormula.coordSwitchingFailureProbability_le_ratio_pow_card, FABL.DNFFormula.toBooleanFunction, FABL.DNFFormula.width, FABL.SwitchingCode, FABL.SwitchingCode.card_le_scale_pow, FABL.SwitchingCode.card_width_one, FABL.coordSwitchingFailureProbability, FABL.coordSwitchingFailureProbability_eq, FABL.coordSwitchingFailureProbability_le_one, FABL.switchingAtomRatio, FABL.switchingAtomRatio_nonneg, FABL.switchingFailureProbability}
\difficulty{6}
Let $\varphi$ be a DNF formula on $n$ Boolean coordinates whose width is at most a
natural number $w$, and let $k$ be a natural number. For every restriction parameter
$\delta\in[0,1]$, the probability that a random restriction leaves the associated
Boolean function with decision-tree depth at least $k$ is at most $(5\delta w)^k$.
\end{theorem}

\begin{theorem}[Hastad switching lemma]
\lean{FABL.hastadSwitchingLemma}
\statementsource{lmn-constant-depth-fourier}{pdf-63759bd909d5-p005-b004}
\statementsource{arora-barak-ch13-circuit-lowerbounds}{pdf-52929dbaa962-p002-b003}
\proofsource{Analysis-of-Boolean-Functions-by-Ryan-ODonnell}{pdf-ffbe9c7b3306-p100-b005}
\label{FABL.hastadSwitchingLemma}
\leanok
\uses{FABL.BooleanFunction, FABL.CNFFormula.booleanDual, FABL.CNFFormula.switchAndOr, FABL.CNFFormula.switchAndOr_toBooleanFunction, FABL.CNFFormula.toBooleanFunction, FABL.CNFFormula.width, FABL.DNFFormula.toBooleanFunction, FABL.DNFFormula.width, FABL.HasCNFWidthLE, FABL.HasDNFWidthLE, FABL.hastadSwitchingLemma_dnf, FABL.switchingFailureProbability, FABL.switchingFailureProbability_booleanDual}
\difficulty{4}
% dataset-pipeline: stub -- prose pending build_informal_statements.py
Statement pending regeneration.
\end{theorem}
